\input{preamble}

\begin{document}

\title{Functional Analysis}
\maketitle
\phantomsection
\label{section-phantom-functional-analysis}

\tableofcontents

\section{Measure spaces}
\label{section-measure-spaces}

\noindent These notes merge the home notes on
measure theory and functional analysis with
the class notes from Bruno Braga's course.
The point is to keep the material in the
style of the stack: definitions, lemmas,
theorems, examples, exercises, and
cross-references by labels.

\begin{definition}
\label{definition-sigma-algebra}
Let $X$ be a set. An {\it algebra of sets} on
$X$ is a nonempty collection $\mathcal A$ of
subsets of $X$ closed under finite unions and
complements. A {\it $\sigma$-algebra} on $X$
is an algebra of sets closed under countable
unions.
\end{definition}

\begin{definition}
\label{definition-generated-sigma-algebra}
Let $\mathcal E \subset \mathcal P(X)$. The
{\it $\sigma$-algebra generated by $\mathcal
E$} is the intersection of all
$\sigma$-algebras containing $\mathcal E$.
\end{definition}

\begin{definition}
\label{definition-borel-algebra}
If $X$ is a topological space, the
$\sigma$-algebra generated by the open sets
of $X$ is called the {\it Borel
$\sigma$-algebra} of $X$ and is denoted
$\mathcal B_X$.
\end{definition}

\begin{definition}
\label{definition-measure}
Let $\mathcal M$ be a $\sigma$-algebra on
$X$. A {\it measure} on $(X,\mathcal M)$ is a
map
$$
\mu : \mathcal M \longrightarrow [0,\infty]
$$
such that $\mu(\emptyset)=0$ and such that
whenever $E_1,E_2,\ldots$ are pairwise
disjoint measurable sets, one has
$$
\mu\left(\bigcup_{n=1}^\infty E_n\right) =
\sum_{n=1}^\infty \mu(E_n).
$$
The triple $(X,\mathcal M,\mu)$ is called a
{\it measure space}.
\end{definition}

\begin{lemma}
\label{lemma-basic-properties-measures}
Let $(X,\mathcal M,\mu)$ be a measure space.
\begin{enumerate}
\item If $E \subset F$, then $\mu(E) \leq
\mu(F)$.
\item If $E_1,E_2,\ldots$ are measurable,
then
$$
\mu\left(\bigcup_{n=1}^\infty E_n\right) \leq
\sum_{n=1}^\infty \mu(E_n).
$$
\item If $E_1\subset E_2\subset\ldots$, then
$$
\mu\left(\bigcup_{n=1}^\infty
E_n\right)=\lim_n \mu(E_n).
$$
\item If $E_1\supset E_2\supset\ldots$ and
$\mu(E_1)<\infty$, then
$$
\mu\left(\bigcap_{n=1}^\infty
E_n\right)=\lim_n\mu(E_n).
$$
\end{enumerate}
\end{lemma}

\begin{definition}
\label{definition-null-set-complete-measure}
A measurable set $N$ is a {\it null set} if
$\mu(N)=0$. A measure is {\it complete} if
every subset of a null set is measurable.
\end{definition}

\begin{theorem}
\label{theorem-completion-measure}
Every measure has a completion. More
precisely, if $(X,\mathcal M,\mu)$ is a
measure space and
$$
\mathcal N = \{N\in \mathcal M : \mu(N)=0\},
$$
then the collection
$$
\overline{\mathcal M} = \{E\cup A :
E\in\mathcal M,\ A\subset N \text{ for some
}N\in\mathcal N\}
$$
is a $\sigma$-algebra, and $\mu$ extends
uniquely to a complete measure on
$\overline{\mathcal M}$.
\end{theorem}

\begin{definition}
\label{definition-outer-measure}
An {\it outer measure} on $X$ is a function
$\mu^*:\mathcal P(X)\to[0,\infty]$ such that
\begin{enumerate}
\item $\mu^*(\emptyset)=0$;
\item $A\subset B$ implies
$\mu^*(A)\leq\mu^*(B)$;
\item $\mu^*(\bigcup A_n)\leq \sum
\mu^*(A_n)$.
\end{enumerate}
\end{definition}

\begin{definition}
\label{definition-carath-measurable}
Let $\mu^*$ be an outer measure on $X$. A set
$A\subset X$ is {\it $\mu^*$-measurable} if
for every $E\subset X$ one has
$$
\mu^*(E)=\mu^*(E\cap A)+\mu^*(E\cap A^c).
$$
\end{definition}

\begin{theorem}[Caratheodory]
\label{theorem-caratheodory}
If $\mu^*$ is an outer measure on $X$, then
the collection of $\mu^*$-measurable subsets
of $X$ is a $\sigma$-algebra and the
restriction of $\mu^*$ to it is a complete
measure.
\end{theorem}

\begin{definition}
\label{definition-premeasure}
Let $\mathcal A$ be an algebra of subsets of
$X$. A function $\mu_0:\mathcal
A\to[0,\infty]$ is a {\it premeasure} if
$\mu_0(\emptyset)=0$ and if countable
additivity holds whenever the countable
disjoint union still belongs to $\mathcal A$.
\end{definition}

\begin{theorem}
\label{theorem-extension-premeasure}
Let $\mu_0$ be a premeasure on an algebra
$\mathcal A$, and let $\mathcal M$ be the
$\sigma$-algebra generated by $\mathcal A$.
Then $\mu_0$ extends to a measure on
$\mathcal M$. If $\mu_0$ is $\sigma$-finite,
this extension is unique.
\end{theorem}

\begin{definition}
\label{definition-borel-measure-r}
A {\it Borel measure on $\mathbb{R}$} is a
measure on $\mathbb{R}$ whose domain is the
Borel $\sigma$-algebra $\mathcal
B_{\mathbb{R}}$.
\end{definition}

\begin{proposition}
\label{proposition-premeasure-half-open}
Let $F:\mathbb{R}\to\mathbb{R}$ be an
increasing and right-continuous function. The
collection $\mathcal A$ of finite disjoint
unions of half-open intervals $(a,b]$
(including $(a,\infty)$ and $(-\infty,b]$) is
an algebra. The function $\mu_0$ defined by
$$
\mu_0\left(\bigcup_{j=1}^n (a_j, b_j]\right)
= \sum_{j=1}^n (F(b_j) - F(a_j))
$$
is a premeasure on $\mathcal A$.
\end{proposition}

\begin{theorem}
\label{theorem-lebesgue-stieltjes}
If $F:\mathbb{R}\to\mathbb{R}$ is an
increasing and right-continuous function,
there is a unique Borel measure $\mu_F$ on
$\mathbb{R}$ such that $\mu_F((a,b]) =
F(b)-F(a)$ for all $a,b\in\mathbb{R}$. The
completion of $\mu_F$ is called the {\it
Lebesgue--Stieltjes measure} associated with
$F$.
\end{theorem}

\begin{example}
\label{example-lebesgue-measure}
The {\it Lebesgue measure} $m$ on
$\mathbb{R}$ is the Lebesgue--Stieltjes
measure associated with the function
$F(x)=x$. It is the unique measure on
$(\mathbb{R}, \mathcal B_{\mathbb{R}})$ that
assigns to each interval its length.
\end{example}

\section
{Measurable functions and integration}
\label
{section-measurable-functions-integration}

\begin{definition}
\label{definition-measurable-map}
Let $(X,\mathcal M)$ and $(Y,\mathcal N)$ be
measurable spaces. A map $f:X\to Y$ is {\it
measurable} if $f^{-1}(E)\in\mathcal M$ for
every $E\in\mathcal N$.
\end{definition}

\begin{lemma}
\label{lemma-measurable-generated}
If $\mathcal N$ is generated by a collection
$\mathcal E$, then $f:X\to Y$ is measurable
if and only if $f^{-1}(E)\in\mathcal M$ for
all $E\in\mathcal E$.
\end{lemma}

\begin{example}
\label{example-continuous-borel-measurable}
If $X,Y$ are topological spaces and $f:X\to
Y$ is continuous, then $f$ is $(\mathcal
B_X,\mathcal B_Y)$-measurable.
\end{example}

\begin{definition}
\label{definition-simple-function}
A complex-valued function $\varphi$ on $X$ is
{\it simple} if it is a finite linear
combination of characteristic functions of
measurable sets. Equivalently, $\varphi$ is
measurable and has finite image.
\end{definition}

\begin{theorem}
\label
{theorem-approximation-simple-functions}
Let $(X,\mathcal M)$ be a measurable space.
\begin{enumerate}
\item If $f:X\to[0,\infty]$ is measurable,
then there is an increasing sequence of
nonnegative simple functions $\varphi_n$ with
$\varphi_n\to f$ pointwise.
\item If $f:X\to\mathbb{C}$ is measurable,
then there is a sequence of simple functions
$\varphi_n$ with $\varphi_n\to f$ pointwise
and $|\varphi_n|\leq |f|$.
\end{enumerate}
\end{theorem}

\begin{proof}
For (1), for each $n\in\mathbb{N}$ and $0\leq
k\leq 2^{2n}-1$, let
$$
E_n^k = f^{-1}((k2^{-n}, (k+1)2^{-n}]) \quad
\text{and} \quad F_n = f^{-1}((2^n, \infty]).
$$
Define
$$
\varphi_n = \sum_{k=0}^{2^{2n}-1}
k2^{-n}\chi_{E_n^k} + 2^n\chi_{F_n}.
$$
One checks that $\varphi_n \leq
\varphi_{n+1}$ and $\varphi_n \to f$
pointwise. For (2), decompose $f = f_1 - f_2
+ i(f_3 - f_4)$ where $f_j \geq 0$ are the
positive and negative parts of the real and
imaginary parts of $f$. Apply part (1) to
each $f_j$.
\end{proof}

\begin{definition}
\label
{definition-integral-nonnegative-simple}
Let $\varphi=\sum_{j=1}^n a_j\chi_{E_j}$ be a
nonnegative simple function written in
standard form. Define
$$
\int \varphi\,d\mu = \sum_{j=1}^n
a_j\mu(E_j),
$$
with the convention $0\cdot\infty=0$.
\end{definition}

\begin{definition}
\label{definition-integral-nonnegative}
For $f:X\to[0,\infty]$ measurable, define
$$
\int f\,d\mu = \sup\left\{\int \varphi\,d\mu:
0\leq \varphi\leq f,\ \varphi\text{
simple}\right\}.
$$
\end{definition}

\begin{theorem}[Monotone convergence]
\label{theorem-monotone-convergence}
If $0\leq f_1\leq f_2\leq\ldots$ are
measurable and $f_n\to f$ pointwise, then
$$
\int f\,d\mu = \lim_n \int f_n\,d\mu.
$$
\end{theorem}

\begin{lemma}[Fatou]
\label{lemma-fatou}
If $f_n\geq 0$ are measurable, then
$$
\int \liminf_n f_n\,d\mu \leq \liminf_n \int
f_n\,d\mu.
$$
\end{lemma}

\begin{theorem}[Dominated convergence]
\label{theorem-dominated-convergence}
Let $f_n$ be measurable functions with
$f_n\to f$ almost everywhere. If there is an
integrable function $g$ such that $|f_n|\leq
g$ for all $n$, then $f$ is integrable and
$$
\lim_n\int |f_n-f|\,d\mu=0.
$$
In particular, $\int f_n\,d\mu\to\int
f\,d\mu$.
\end{theorem}

\begin{definition}
\label{definition-lp-space}
Let $1\leq p<\infty$. The space $L^p(X,\mu)$
is the vector space of measurable functions
$f$ with
$$
\|f\|_p = \left(\int
|f|^p\,d\mu\right)^{1/p}<\infty,
$$
modulo equality almost everywhere. Also
$L^\infty(X,\mu)$ is the space of essentially
bounded measurable functions with the
essential supremum norm.
\end{definition}

\begin{theorem}[Holder inequality]
\label{theorem-holder-inequality}
Let $1<p,q<\infty$ and $1/p+1/q=1$. If $f\in
L^p$ and $g\in L^q$, then $fg\in L^1$ and
$$
\|fg\|_1\leq \|f\|_p\|g\|_q.
$$
\end{theorem}

\begin{theorem}[Minkowski inequality]
\label{theorem-minkowski-inequality}
For $1\leq p\leq\infty$, $L^p(X,\mu)$ is a
normed vector space. Equivalently, if $f,g\in
L^p$, then
$$
\|f+g\|_p\leq \|f\|_p+\|g\|_p.
$$
\end{theorem}

\begin{theorem}
\label{theorem-lp-banach}
For $1\leq p\leq\infty$, the space
$L^p(X,\mu)$ is a Banach space.
\end{theorem}

\begin{example}
\label{example-sequence-lp}
For $1\leq p<\infty$,
$$
\ell^p = \left\{x=(x_n):\sum_{n=1}^\infty
|x_n|^p<\infty\right\}
$$
with $\|x\|_p=(\sum |x_n|^p)^{1/p}$ is a
Banach space. Also $\ell^\infty$ is a Banach
space with $\|x\|_\infty=\sup_n |x_n|$, and
$c_0\subset\ell^\infty$ is the closed
subspace of sequences converging to zero.
\end{example}

\begin{definition}
\label{definition-convergence-in-measure}
Let $(X,\mathcal M,\mu)$ be a measure space.
A sequence $f_n$ converges to $f$ {\it in
measure} if for every $\varepsilon>0$,
$$
\mu\{x: |f_n(x)-f(x)|>\varepsilon\}\to 0.
$$
\end{definition}

\begin{theorem}
\label
{theorem-l1-convergence-implies-measure}
If $f_n\to f$ in $L^1$, then $f_n\to f$ in
measure. Moreover, some subsequence converges
to $f$ almost everywhere.
\end{theorem}

\begin{theorem}[Egoroff]
\label{theorem-egoroff}
Assume $\mu(X)<\infty$. If $f_n\to f$ almost
everywhere, then for every $\varepsilon>0$
there is a measurable set $E$ with
$\mu(E)<\varepsilon$ such that $f_n\to f$
uniformly on $X\setminus E$.
\end{theorem}

\begin{definition}
\label{definition-almost-uniform-convergence}
A sequence of measurable functions $f_n$ is
said to converge to $f$ {\it almost
uniformly} if for every $\varepsilon>0$ there
exists a measurable set $E$ with
$\mu(E)<\varepsilon$ such that $f_n \to f$
uniformly on $E^c$.
\end{definition}

\begin{remark}
\label{remark-egoroff-almost-uniform}
Egoroff's theorem can be restated as: on a
finite measure space, almost everywhere
convergence implies almost uniform
convergence.
\end{remark}

\section{Product measures}
\label{section-product-measures}

\begin{definition}
\label{definition-product-sigma-algebra}
Let $(X,\mathcal M)$ and $(Y,\mathcal N)$ be
measurable spaces. The {\it product
$\sigma$-algebra} $\mathcal M\otimes\mathcal
N$ is the $\sigma$-algebra on $X\times Y$
generated by the measurable rectangles
$A\times B$, where $A\in\mathcal M$ and
$B\in\mathcal N$.
\end{definition}

\begin{proposition}
\label{proposition-product-borel}
Let $X_1, \ldots, X_n$ be metric spaces and
$X = \prod X_i$ with the product metric. Then
$\bigotimes \mathcal B_{X_i} \subset \mathcal
B_X$. If each $X_i$ is separable, then
$\bigotimes \mathcal B_{X_i} = \mathcal B_X$.
\end{proposition}

\begin{theorem}
\label{theorem-product-measure}
Let $(X,\mathcal M,\mu)$ and $(Y,\mathcal
N,\nu)$ be $\sigma$-finite measure spaces.
There is a unique measure $\mu\times\nu$ on
$\mathcal M\otimes\mathcal N$ such that
$$
(\mu\times\nu)(A\times B)=\mu(A)\nu(B)
$$
for all measurable rectangles $A\times B$.
\end{theorem}

\begin{definition}
\label{definition-sections}
If $E\subset X\times Y$, define
$$
E_x = \{y\in Y:(x,y)\in E\}, \quad E^y =
\{x\in X:(x,y)\in E\}.
$$
If $f:X\times Y\to\mathbb{C}$, define
$f_x(y)=f^y(x)=f(x,y)$.
\end{definition}

\begin{theorem}[Tonelli]
\label{theorem-tonelli}
Let $(X,\mathcal M,\mu)$ and $(Y,\mathcal
N,\nu)$ be $\sigma$-finite measure spaces. If
$f:X\times Y\to[0,\infty]$ is measurable,
then $x\mapsto\int f_x\,d\nu$ and
$y\mapsto\int f^y\,d\mu$ are measurable, and
$$
\int_{X\times Y} f\,d(\mu\times\nu) =
\int_X\left(\int_Y
f(x,y)\,d\nu(y)\right)d\mu(x) =
\int_Y\left(\int_X
f(x,y)\,d\mu(x)\right)d\nu(y).
$$
\end{theorem}

\begin{theorem}[Fubini]
\label{theorem-fubini}
Under the hypotheses of Theorem
\ref{theorem-tonelli}, if $f\in L^1(X\times
Y)$, then almost all sections are integrable
and the same iterated-integral formula holds.
\end{theorem}

\section{Compactness and spaces of functions}
\label{section-compactness-function-spaces}

\begin{definition}
\label{definition-uniform-norm}
Let $X$ be a set. The space $B(X)$ of bounded
scalar-valued functions on $X$ is normed by
$$
\|f\|_u = \sup_{x\in X} |f(x)|.
$$
\end{definition}

\begin{proposition}
\label{proposition-bx-banach}
The space $B(X)$ is a Banach space. If $X$ is
a topological space, then $BC(X)$, the space
of bounded continuous functions, is a closed
subspace of $B(X)$ and is therefore Banach.
\end{proposition}

\begin{definition}
\label{definition-compact}
A topological space $X$ is {\it compact} if
every open cover of $X$ has a finite
subcover. A subset is {\it precompact} if its
closure is compact.
\end{definition}

\begin{lemma}
\label{lemma-compact-closed-fip}
A topological space $X$ is compact if and
only if every family of closed sets with the
finite intersection property has nonempty
intersection.
\end{lemma}

\begin{theorem}
\label
{theorem-compact-metric-characterization}
Let $(X,d)$ be a metric space and $E\subset
X$. Then the following are equivalent.
\begin{enumerate}
\item $E$ is compact.
\item Every sequence in $E$ has a subsequence
converging to a point of $E$.
\item $E$ is complete and totally bounded.
\end{enumerate}
\end{theorem}

\begin{theorem}[Urysohn]
\label{theorem-urysohn}
Let $X$ be a normal topological space. If $A$
and $B$ are disjoint closed subsets of $X$,
then there is a continuous function
$f:X\to[0,1]$ such that $f|_A=0$ and
$f|_B=1$.
\end{theorem}

\begin{theorem}[Tietze]
\label{theorem-tietze}
Let $X$ be a normal topological space, let
$A\subset X$ be closed, and let $f:A\to[a,b]$
be continuous. Then there is a continuous
$F:X\to[a,b]$ such that $F|_A=f$.
\end{theorem}

\begin{definition}
\label{definition-equicontinuous}
Let $X$ be a topological space. A family
$\mathcal F \subset C(X)$ is {\it
equicontinuous} if for every $x\in X$ and
every $\varepsilon>0$, there is a
neighborhood $U$ of $x$ such that
$|f(x)-f(y)|<\varepsilon$ for all $y\in U$
and all $f\in\mathcal F$.
\end{definition}

\begin{theorem}[Arzelà--Ascoli]
\label{theorem-arzela-ascoli}
Let $X$ be a compact Hausdorff space. A
subset $\mathcal F \subset C(X)$ is
precompact in the uniform norm if and only if
it is equicontinuous and pointwise bounded.
\end{theorem}

\begin{theorem}[Stone--Weierstrass]
\label{theorem-stone-weierstrass}
Let $X$ be a compact Hausdorff space. If
$\mathcal A$ is a subalgebra of $C(X,
\mathbb{R})$ that separates points and
contains the constant functions, then
$\mathcal A$ is dense in $C(X, \mathbb{R})$
for the uniform norm. If $\mathcal A$ is a
subalgebra of $C(X, \mathbb{C})$ that
separates points, contains the constant
functions, and is closed under complex
conjugation, then $\mathcal A$ is dense in
$C(X, \mathbb{C})$.
\end{theorem}

\begin{definition}
\label{definition-cc-c0}
Let $X$ be a locally compact Hausdorff space.
\begin{enumerate}
\item $C_c(X)$ is the space of continuous
functions with compact support.
\item $C_0(X)$ is the space of continuous
functions that vanish at infinity, i.e., for
every $\varepsilon>0$, the set $\{x\in X :
|f(x)|\geq\varepsilon\}$ is compact.
\end{enumerate}
\end{definition}

\begin{proposition}
\label{proposition-c0-completion-cc}
If $X$ is a locally compact Hausdorff space,
$C_0(X)$ is the completion of $C_c(X)$ in the
uniform norm.
\end{proposition}

\section{Normed spaces and Banach spaces}
\label{section-normed-banach-spaces}

\begin{definition}
\label{definition-norm}
Let $X$ be a vector space over a field
$\mathbb{K} \in \{\mathbb{R}, \mathbb{C}\}$.
A function $\|\cdot\|: X \to [0, \infty)$ is
a {\it norm} if:
\begin{enumerate}
\item $\|x\| = 0$ if and only if $x = 0$;
\item $\|\lambda x\| = |\lambda|\|x\|$ for
all $\lambda \in \mathbb{K}$ and $x \in X$;
\item $\|x+y\| \leq \|x\| + \|y\|$ for all
$x, y \in X$.
\end{enumerate}
The pair $(X, \|\cdot\|)$ is called a {\it
normed space}. A normed space is a metric
space with the metric $d(x,y) = \|x-y\|$. If
$(X, \|\cdot\|)$ is complete as a metric
space, it is called a {\it Banach space}.
\end{definition}

\begin{example}
\label{example-banach-spaces}
The following are important examples of
Banach spaces:
\begin{enumerate}
\item {\bf Sequence spaces.} For $1 \leq p <
\infty$, the space
$$
\ell_p = \left\{ (x_n) \in
\mathbb{K}^\mathbb{N} : \sum_{n=1}^\infty
|x_n|^p < \infty \right\}
$$
with the norm $\|(x_n)\|_p = (\sum
|x_n|^p)^{1/p}$ is a Banach space.
\item {\bf Bounded sequences.} The space
$$
\ell_\infty = \{ (x_n) \in
\mathbb{K}^\mathbb{N} : \sup_n |x_n| < \infty
\}
$$
with the norm $\|(x_n)\|_\infty = \sup_n
|x_n|$ is a Banach space.
\item {\bf Null sequences.} The space $c_0 =
\{ (x_n) \in \ell_\infty : \lim_n x_n = 0 \}$
is a closed subspace of $\ell_\infty$ and
thus a Banach space.
\item {\bf Continuous functions.} If $K$ is a
compact Hausdorff space, $C(K)$ with the
supremum norm $\|f\| = \sup_{x \in K} |f(x)|$
is a Banach space.
\item {\bf $L^p$ spaces.} Let $(X,
\mathcal{A}, \mu)$ be a measure space. For $1
\leq p < \infty$, the space $L^p(\mu)$ of
equivalence classes of measurable functions
such that $\int |f|^p d\mu < \infty$ is a
Banach space.
\end{enumerate}
\end{example}

\begin{lemma}[Young's inequality]
\label{lemma-young-inequality}
Let $a, b \geq 0$ and $p, q > 1$ such that
$1/p + 1/q = 1$. Then
$$
ab \leq \frac{a^p}{p} + \frac{b^q}{q}.
$$
\end{lemma}

\begin{proof}
Consider the function $f(x) = x^{p-1}$. The
area under $f$ from $0$ to $a$ is $\int_0^a
x^{p-1} dx = a^p/p$. The inverse function is
$g(y) = y^{1/(p-1)} = y^{q-1}$. The area
under $g$ from $0$ to $b$ is $\int_0^b
y^{q-1} dy = b^q/q$. By drawing a figure, it
is clear that for any $a, b > 0$, the sum of
these areas is greater than or equal to the
area of the rectangle with sides $a$ and $b$.
\end{proof}

\begin{theorem}[Hölder's inequality]
\label{theorem-holder-inequality}
Let $p, q \in (1, \infty)$ with $1/p + 1/q =
1$. If $(x_n) \in \ell_p$ and $(y_n) \in
\ell_q$, then $(x_n y_n) \in \ell_1$ and
$$
\|(x_n y_n)\|_1 \leq \|(x_n)\|_p \|(y_n)\|_q.
$$
\end{theorem}

\begin{proof}
If $\|(x_n)\|_p = 0$ or $\|(y_n)\|_q = 0$,
the inequality is trivial. Assume
$\|(x_n)\|_p = \|(y_n)\|_q = 1$. By Young's
inequality, for each $n$:
$$
|x_n y_n| \leq \frac{|x_n|^p}{p} +
\frac{|y_n|^q}{q}.
$$
Summing over $n$:
$$
\sum |x_n y_n| \leq \frac{1}{p} \sum |x_n|^p
+ \frac{1}{q} \sum |y_n|^q = \frac{1}{p} +
\frac{1}{q} = 1.
$$
For the general case, normalize $\tilde{x}_n
= x_n / \|x\|_p$ and $\tilde{y}_n = y_n /
\|y\|_q$.
\end{proof}

\begin{theorem}[Minkowski's inequality]
\label{theorem-minkowski-inequality}
If $1 \leq p < \infty$ and $x, y \in \ell_p$,
then $x+y \in \ell_p$ and $\|x+y\|_p \leq
\|x\|_p + \|y\|_p$.
\end{theorem}

\begin{proof}
For $p=1$ it is the triangle inequality. For
$p > 1$:
$$
\begin{aligned}
\|x+y\|_p^p &= \sum |x_n+y_n|^p \leq \sum
|x_n||x_n+y_n|^{p-1} + \sum
|y_n||x_n+y_n|^{p-1} \\ &\leq \|x\|_p
\|(|x_n+y_n|^{p-1})\|_q + \|y\|_p
\|(|x_n+y_n|^{p-1})\|_q \\
&\quad \text{(Hölder)}
\end{aligned}
$$

\noindent
Note that $(p-1)q = p$. Thus
$\|(|x_n+y_n|^{p-1})\|_q = (\sum
|x_n+y_n|^p)^{1/q} = \|x+y\|_p^{p/q}$.
Dividing by $\|x+y\|_p^{p/q}$ (assuming $x+y
\neq 0$) and using $p - p/q = 1$, we get the
result.
\end{proof}

\begin{theorem}
\label{theorem-lp-completeness}
The space $\ell_p$ is a Banach space for $1
\leq p \leq \infty$.
\end{theorem}

\begin{proof}
We treat the case $1 \leq p < \infty$. Let
$(x^{(n)})$ be a Cauchy sequence in $\ell_p$,
where $x^{(n)} = (x_k^{(n)})_k$. For each
fixed $k$, $|x_k^{(n)} - x_k^{(m)}| \leq
\|x^{(n)} - x^{(m)}\|_p$, so $(x_k^{(n)})_n$
is a Cauchy sequence in $\mathbb{K}$. Let
$x_k = \lim_n x_k^{(n)}$ and $x = (x_k)_k$.
Since $(x^{(n)})$ is Cauchy, it is bounded:
$\|x^{(n)}\|_p \leq M$ for all $n$. For any
$N \in \mathbb{N}$:
$$
\sum_{k=1}^N |x_k^{(n)}|^p \leq M^p.
$$
Taking $n \to \infty$, $\sum_{k=1}^N |x_k|^p
\leq M^p$. As $N$ is arbitrary, $x \in
\ell_p$. Now show $x^{(n)} \to x$ in norm.
Given $\varepsilon > 0$, $\exists n_0$ s.t.
$m, n > n_0 \implies \|x^{(n)} - x^{(m)}\|_p
< \varepsilon$. For any $N$:
$$
\sum_{k=1}^N |x_k^{(n)} - x_k^{(m)}|^p <
\varepsilon^p.
$$
Taking $m \to \infty$: $\sum_{k=1}^N
|x_k^{(n)} - x_k|^p \leq \varepsilon^p$. As
$N$ is arbitrary, $\|x^{(n)} - x\|_p \leq
\varepsilon$.
\end{proof}

\begin{remark}
\label{remark-lp-inclusions}
\begin{enumerate}
\item If $p \leq q$, then $\ell_p \subseteq
\ell_q$.
\item If $\mu(X) < \infty$ and $p \leq q$,
then $L^q(\mu) \subseteq L^p(\mu)$.
\item The space $c_{00}$ of sequences with
finite support is dense in $\ell_p$ for $p <
\infty$, but not in $\ell_\infty$.
\end{enumerate}
\end{remark}



\subsection{Equivalent norms, embeddings, and duals}
\label{subsection-equivalent-norms-embeddings-duals}

\begin{definition}
\label{definition-equivalent-norms}
Let $X$ be a vector space and let
$\|\cdot\|$ and $\|\cdot\|'$ be norms on
$X$. They are {\it equivalent} if there is
$L>0$ such that
$
\frac1L\|x\|'\leq \|x\|\leq
L\|x\|'
\qquad\text{for every }x\in X.
$
\end{definition}

\begin{remark}
\label{remark-finite-dimensional-norms-equivalent}
On a finite-dimensional vector space all
norms are equivalent.
\end{remark}

\begin{definition}
\label{definition-isomorphism-isomorphic-embedding}
Let $X,Y$ be normed spaces.
\begin{enumerate}
\item A bijective linear map $T:X\to Y$ is an
{\it isomorphism} if $T$ and $T^{-1}$ are
continuous.
\item A linear map $T:X\to Y$ is an
{\it isomorphic embedding} if
$T:X\to T(X)$ is an isomorphism.
\end{enumerate}
Equivalently, $T$ is an isomorphic embedding
if there exists $L>0$ such that
$
\frac1L\|x\|_X\leq\|Tx\|_Y\leq L\|x\|_X
\qquad\text{for every }x\in X.
$
\end{definition}

\begin{proof}[Proof of Lemma \ref{lemma-riesz}]
Choose $x\in X\setminus Y$. Since
$d(x,Y)/a>d(x,Y)$, there exists $y\in Y$ such
that
$
\|x-y\|\leq \frac{d(x,Y)}a.
$
Put
$
z=\frac{x-y}{\|x-y\|}.
$
Then $\|z\|=1$. For every $w\in Y$,
\begin{align*}
\|w-z\|
&=
\frac1{\|x-y\|}
\left\|\|x-y\|w+y-x\right\|\\
&\geq \frac{d(x,Y)}{\|x-y\|}\\
&\geq a,
\end{align*}
because $\|x-y\|w+y\in Y$. Hence
$d(z,Y)\geq a$.
\end{proof}

\begin{proof}
[Proof of Theorem
\ref{theorem-finite-dimensional-operators-%
continuous}]
We argue by induction on $\dim X$. The
zero-dimensional case is immediate. Assume
the result for dimension $n$ and let
$\dim X=n+1$. Choose an $n$-dimensional
subspace $Y\subset X$ and $x\in X\setminus
Y$. By the induction hypothesis, linear maps
restricted to $Y$ are continuous. Writing
each vector uniquely as $y+\lambda x$, the
problem reduces to controlling the scalar
coefficient $\lambda$. Riesz's lemma gives a
positive lower bound for the distance from
the unit vector in the $x$-direction to $Y$,
and therefore bounds $|\lambda|$ by a
constant multiple of $\|y+\lambda x\|$.
Consequently every linear map on $X$ is
bounded.

The Portuguese course notes present this
inductive argument but mark the corresponding
proof that finite-dimensional spaces are
Banach as incomplete.
\end{proof}

\begin{proof}
[Proof of Theorem
\ref{theorem-unit-ball-compact-finite-dim}]
If $\dim X=n$, choose an algebraic
isomorphism $T:\mathbb K^n\to X$. By
Theorem
\ref{theorem-finite-dimensional-operators-%
continuous},
$T$ and $T^{-1}$ are continuous. Hence
$B_X$ is homeomorphic to a closed bounded
subset of $\mathbb K^n$, and is compact.

Conversely, suppose $X$ is infinite
dimensional. By Riesz's lemma, inductively
choose $x_n\in S_X$ such that
$
d\left(x_n,
\operatorname{span}\{x_1,\ldots,x_{n-1}\}
\right)>\frac12.
$
Then $\|x_n-x_m\|>1/2$ whenever $n\neq m$.
Thus $(x_n)$ has no convergent subsequence,
so $B_X$ is not compact.
\end{proof}

\begin{definition}
\label{definition-topological-dual}
The {\it topological dual} of a normed space
$X$ is
$
X^*=\mathcal L(X,\mathbb K).
$
\end{definition}

\begin{theorem}
\label{theorem-dual-lp-lq-course}
Let $1<p,q<\infty$ and
$1/p+1/q=1$. Then
$
\ell_p^*\cong \ell_q
$
isometrically.
\end{theorem}

\begin{proof}
For $y=(y_n)\in\ell_q$, define
$
\phi(y)(x)=\sum_{n=1}^\infty x_ny_n,
\qquad x=(x_n)\in\ell_p.
$
Hölder's inequality shows that the series is
absolutely convergent and
$
\|\phi(y)\|\leq\|y\|_q.
$
To obtain the reverse inequality, for
$y\neq0$ choose
$
x_n=
\operatorname{sign}(y_n)
\frac{|y_n|^{q-1}}
{\|y\|_q^{q/p}}.
$
Since $(q-1)p=q$, one has $\|x\|_p=1$, and
$
\phi(y)(x)=\|y\|_q.
$
Thus $\phi$ is an isometry.

To prove surjectivity, let $f\in\ell_p^*$ and
put $y_n=f(e_n)$. For each $k$, set
$
x_n=
\begin{cases}
\operatorname{sign}(y_n)|y_n|^{q-1},
&n\leq k,\\
0,&n>k.
\end{cases}
$
Then
\begin{align*}
\sum_{n=1}^k|y_n|^q
&=
f\left(\sum_{n=1}^k x_ne_n\right)\\
&\leq
\|f\|
\left(\sum_{n=1}^k |y_n|^{(q-1)p}\right)^{1/p}\\
&=
\|f\|
\left(\sum_{n=1}^k |y_n|^q\right)^{1/p}.
\end{align*}
Hence the partial $\ell_q$-norms of $y$ are
bounded by $\|f\|$, so $y\in\ell_q$.
On finitely supported vectors,
$f=\phi(y)$, and density of $c_{00}$ in
$\ell_p$ gives equality on all of $\ell_p$.
\end{proof}

\begin{example}
\label{example-duals-c0-l1-linfty}
The course notes also record
$
c_0^*\cong\ell_1,\qquad
\ell_1^*\cong\ell_\infty,
$
and point out that the analogous naive
argument does not give
$\ell_\infty^*\cong\ell_1$.
\end{example}

\begin{definition}
\label{definition-weak-convergence-sequence}
A sequence $(x_n)$ in a normed space $X$
{\it converges weakly} to $x\in X$ if
$
f(x_n)\longrightarrow f(x)
\qquad\text{for every }f\in X^*.
$
We write $x_n\rightharpoonup x$.
\end{definition}

\begin{example}
\label{example-en-weak-lp}
For $1<p<\infty$, the canonical sequence
$(e_n)$ in $\ell_p$ converges weakly to zero.
Indeed, every $f\in\ell_p^*$ is represented
by some $y\in\ell_q$, and
$f(e_n)=y_n\to0$.
\end{example}

\begin{proposition}
\label{proposition-bounded-operators-preserve-weak-convergence}
If $T:X\to Y$ is bounded and
$x_n\rightharpoonup x$ in $X$, then
$Tx_n\rightharpoonup Tx$ in $Y$.
\end{proposition}

\begin{proof}
Let $f\in Y^*$. Then $T^*f=f\circ T\in X^*$,
and therefore
$
f(Tx_n)=(T^*f)(x_n)\longrightarrow
(T^*f)(x)=f(Tx).
$
\end{proof}

\begin{theorem}[Pitt, direction proved in the course notes]
\label{theorem-pitt-lp-no-embedding}
Let $1<q<p<\infty$. There is no isomorphic
embedding
$\ell_p\hookrightarrow\ell_q$.
\end{theorem}

\begin{remark}
\label{remark-pitt-source-overstatement}
The Portuguese notes state this theorem with
the broader hypothesis $p\neq q$, but the
proof immediately assumes $p>q$, and a later
remark explicitly says that this proof does
not treat the opposite direction. We record
here only the direction actually proved in
the source.
\end{remark}

\begin{proof}[Proof sketch from the course notes]
The notes treat the case $p>q$. Suppose
$T:\ell_p\to\ell_q$ is an isomorphic
embedding. Since $e_n\rightharpoonup0$ in
$\ell_p$, one has $T(e_n)\rightharpoonup0$
in $\ell_q$. Passing to a subsequence and
using a gliding-hump argument, the notes say
that one may arrange the vectors $T(e_n)$ to
have essentially disjoint supports. The
epsilon details of this choice are explicitly
left to the reader in the source.

For such a subsequence there is $L\geq1$
such that
$
\left\|\sum_{n=1}^kT(e_n)\right\|_q
\geq
\frac1L
\left(\sum_{n=1}^k\|T(e_n)\|_q^q\right)^{1/q}
\geq \frac aL k^{1/q},
$
where
$a=\inf_n\|T(e_n)\|_q>0$ because $T$ is an
isomorphic embedding. On the other hand,
$
\left\|\sum_{n=1}^kT(e_n)\right\|_q
=
\left\|T\left(\sum_{n=1}^ke_n\right)\right\|_q
\leq
\|T\|k^{1/p}.
$
Thus
$
\frac aL k^{1/q}\leq\|T\|k^{1/p}
\qquad\text{for all }k,
$
which is impossible when $p>q$.
\end{proof}

\subsection{Completions}
\label{subsection-completions}

\begin{definition}
\label{definition-completion-normed-space}
A Banach space $\widetilde X$ is a
{\it completion} of a normed space $X$ if
$X$ is isometrically embedded as a dense
subspace of $\widetilde X$.
\end{definition}

\begin{theorem}
\label{theorem-completion-existence-uniqueness}
Every normed space has a completion, unique
up to an isometric isomorphism fixing the
original space.
\end{theorem}

\begin{proof}
Let
$
\widetilde{\widetilde X}
=
\{(x_n)\in X^{\mathbb N}:(x_n)
\text{ is Cauchy}\}.
$
Define
$
(x_n)\sim(y_n)
\quad\Longleftrightarrow\quad
\lim_n\|x_n-y_n\|=0
$
and put
$
\widetilde X=
\widetilde{\widetilde X}/\sim.
$
Vector operations are defined termwise and
$
\left\|[(x_n)]\right\|
=
\lim_n\|x_n\|.
$
The map
$
X\longrightarrow\widetilde X,
\qquad
x\longmapsto[(x,x,\ldots)]
$
is a linear isometry.

To see density, fix $[(x_n)]\in\widetilde X$.
Since $(x_n)$ is Cauchy, the constant
sequences $[(x_m,x_m,\ldots)]$ converge to
$[(x_n)]$.

For completeness, let
$([(x_n^k)])_k$ be Cauchy in $\widetilde X$.
After passing to a subsequence, choose
$x^k\in X$ with
$
\left\|[(x_n^m)]-[(x^k,x^k,\ldots)]\right\|
<\frac1k
\qquad(m\geq k).
$
Then $(x^k)$ is Cauchy in $X$, since for
$m\geq k$,
$
\|x^k-x^m\|
\leq\frac1k+\frac1m.
$
Hence $[(x^k)]\in\widetilde X$, and the
original Cauchy sequence converges to this
class. The source leaves uniqueness as an
exercise.
\end{proof}

\begin{proposition}
\label{proposition-hamel-basis-banach-uncountable}
An infinite-dimensional Banach space cannot
have a countable Hamel basis.
\end{proposition}

\section{Quotient spaces}
\label{section-quotient-spaces}

Given a normed space $X$ and a linear
subspace $Y$, we can consider the quotient
vector space $X/Y$ consisting of cosets $[x]
= x + Y$.

\begin{definition}
\label{definition-quotient-norm}
Let $X$ be a normed space and $Y \subseteq X$
be a subspace. The {\it quotient seminorm} on
$X/Y$ is defined by
$$
\|x+Y\|_{X/Y} = \inf \{ \|x-y\| : y \in Y \}
= d(x, Y).
$$
\end{definition}

\begin{proposition}
\label{proposition-quotient-norm-properties}
Let $X$ be a normed space and $Y$ a subspace.
\begin{enumerate}
\item The function $\|\cdot\|_{X/Y}$ is a
seminorm on $X/Y$.
\item It is a norm if and only if $Y$ is
closed.
\item If $X$ is a Banach space and $Y$ is
closed, then $X/Y$ is a Banach space.
\end{enumerate}
\end{proposition}

\begin{proof}
(1) and (2) are straightforward. For (3),
assume $X$ is Banach and $Y$ is closed. Let
$([x_n])$ be a Cauchy sequence in $X/Y$.
Passing to a subsequence, we may assume
$\|[x_n] - [x_{n+1}]\| < 2^{-n}$. We find a
sequence $(y_n)$ in $X$ such that $y_n \in
[x_n]$ and $\|y_n - y_{n+1}\| < 2^{-n}$ by
choosing $w_n \in Y$ such that $\|x_n -
x_{n+1} - w_n\| < 2^{-n}$ and setting $y_n$
appropriately. Then $(y_n)$ is Cauchy in $X$,
hence converges to some $y \in X$. Then
$[x_n] \to [y]$ in $X/Y$.
\end{proof}

\begin{definition}
\label{definition-canonical-projection}
The map $Q : X \to X/Y$ defined by $Q(x) =
x+Y$ is called the {\it canonical
projection}.
\end{definition}

\begin{proposition}
\label{proposition-canonical-projection-norm}
The canonical projection $Q : X \to X/Y$ is a
linear operator with $\|Q\| \leq 1$. If $Y$
is a proper closed subspace, then $\|Q\| =
1$.
\end{proposition}

\begin{theorem}[First isomorphism theorem]
\label{theorem-isomorphism-theorem-operators}
Let $T : X \to Y$ be a bounded linear
operator. There exists a unique bounded
linear operator $\widehat{T} : X/\ker T \to
Y$ such that $T = \widehat{T} \circ Q$.
Moreover, $\|\widehat{T}\| = \|T\|$.
\end{theorem}

\begin{proof}
The algebraic existence of $\widehat{T}$ is
standard. For the norm:
$$
\begin{aligned}
\|\widehat{T}\| &= \sup \{ \|T(x)\| : x+\ker
T \in B_{X/\ker T} \} \\ &= \sup \{
\|T(x-y)\| : y \in \ker T, \|x+\ker T\| \leq
1 \}.
\end{aligned}
$$

\noindent
Given $\varepsilon > 0$, for each $x$ with
$\|x+\ker T\| \leq 1$, we can find $y \in
\ker T$ such that $\|x-y\| < 1+\varepsilon$.
Thus $\|T(x-y)\| \leq \|T\|(1+\varepsilon)$.
This shows $\|\widehat{T}\| \leq \|T\|$. The
reverse inequality $\|T\| \leq
\|\widehat{T}\|\|Q\| \leq \|\widehat{T}\|$ is
immediate.
\end{proof}


\begin{theorem}
\label{theorem-separable-banach-quotient-l1}
Every separable Banach space $X$ is
isometrically isomorphic to a quotient
$\ell_1/Y$ for some closed subspace
$Y\subseteq\ell_1$.
\end{theorem}

\begin{proof}
Choose a dense sequence $(x_n)$ in $B_X$ and
define
$
T:\ell_1\longrightarrow X,
\qquad
T((\lambda_n))=
\sum_{n=1}^\infty\lambda_nx_n.
$
The series converges absolutely and
$\|T\|\leq1$.

Fix $x\in B_X$ and $\varepsilon\in(0,1)$.
Using density, choose $y_1\in B_{\ell_1}$
with
$
\|x-Ty_1\|<\varepsilon.
$
Inductively choose $y_n$ so that
$
\|y_n\|\leq\varepsilon^{n-1},
\qquad
\left\|x-
T(y_1+\cdots+y_n)\right\|<\varepsilon^n.
$
Then $y=\sum_ny_n$ exists in $\ell_1$ and,
by continuity of $T$, $Ty=x$. Thus $T$ is
surjective.

The first isomorphism theorem gives a
bijection
$
T':\ell_1/\ker T\longrightarrow X.
$
The estimates above show that every
$x\in B_X$ has preimages with norm at most
$1/(1-\varepsilon)$. Letting
$\varepsilon\downarrow0$ yields
$\|(T')^{-1}\|\leq1$, while
$\|T'\|=\|T\|\leq1$. Therefore $T'$ is an
isometry.
\end{proof}

\begin{remark}
\label{remark-linfty-not-separable}
The space $\ell_\infty$ is not separable.
Indeed, for every $A\subseteq\mathbb N$ let
$\chi_A\in\ell_\infty$ be its characteristic
sequence. If $A\neq B$, then
$\|\chi_A-\chi_B\|_\infty=1$, so the balls
$B(\chi_A,1/2)$ form an uncountable pairwise
disjoint family of open sets.
\end{remark}

\section{Linear maps and bounded operators}
\label{section-bounded-operators}

\begin{definition}
\label{definition-bounded-linear-map}
Let $X, Y$ be normed spaces. A linear map $T
: X \to Y$ is {\it bounded} if there is $C
\geq 0$ such that $\|Tx\| \leq C\|x\|$ for
all $x \in X$.
\end{definition}

\begin{proposition}
\label{proposition-bounded-iff-continuous}
Let $T : X \to Y$ be linear. The following
are equivalent:
\begin{enumerate}
\item $T$ is continuous at $0$.
\item $T$ is continuous.
\item $T$ is bounded.
\item $T$ is uniformly continuous.
\end{enumerate}
\end{proposition}

\begin{proof}
The equivalence $(1) \iff (2) \iff (4)$ is
standard for linear maps. For $(2) \implies
(3)$, if $T$ is continuous at $0$, then for
$\varepsilon=1$ there exists $\delta>0$ such
that $\|x\| < \delta \implies \|Tx\| < 1$.
For any $x \neq 0$, $z =
\frac{\delta}{2\|x\|}x$ has norm $\delta/2 <
\delta$, so $\|Tz\| < 1$, which means
$\frac{\delta}{2\|x\|}\|Tx\| < 1$, i.e.,
$\|Tx\| \leq \frac{2}{\delta}\|x\|$.
\end{proof}

\begin{definition}
\label{definition-operator-norm}
The vector space of bounded linear maps from
$X$ to $Y$ is denoted $\mathcal{L}(X, Y)$.
The {\it operator norm} on $\mathcal{L}(X,
Y)$ is defined by
$$
\|T\| = \sup_{x \in B_X} \|Tx\| =
\sup_{\|x\|=1} \|Tx\| = \inf \{ C \geq 0 :
\|Tx\| \leq C\|x\| \text{ for all } x \in X
\}.
$$
\end{definition}

\begin{proposition}
\label{proposition-operator-space-banach}
If $Y$ is a Banach space, then
$\mathcal{L}(X, Y)$ is a Banach space.
\end{proposition}

\begin{proof}
Let $(T_n)$ be a Cauchy sequence in
$\mathcal{L}(X, Y)$. For each $x \in X$,
$\|T_n x - T_m x\| \leq \|T_n - T_m\| \|x\|$,
so $(T_n x)$ is Cauchy in $Y$. Since $Y$ is
complete, let $Tx = \lim_n T_n x$. Linearity
of $T$ is easy. Since $\|T_n\|$ is bounded
(say by $M$), $\|Tx\| = \lim \|T_n x\| \leq
M\|x\|$, so $T$ is bounded. To show $\|T_n -
T\| \to 0$, for fixed $\varepsilon > 0$,
choose $N$ such that $n, m \geq N \implies
\|T_n - T_m\| < \varepsilon$. For $\|x\| \leq
1$, $\|T_n x - T_m x\| < \varepsilon$.
Letting $m \to \infty$, $\|T_n x - Tx\| \leq
\varepsilon$, so $\|T_n - T\| \leq
\varepsilon$.
\end{proof}

\begin{lemma}[Riesz]
\label{lemma-riesz}
Let $X$ be a normed space, $Y \subset X$ a
proper closed subspace, and $a \in (0, 1)$.
Then there exists $x \in X$ with $\|x\|=1$
such that $d(x, Y) \geq a$.
\end{lemma}

\begin{theorem}
\label
{theorem-finite-dimensional-operators-%
continuous}
If $\dim X < \infty$, then every linear map
$T : X \to Y$ is bounded.
\end{theorem}

\begin{theorem}
\label{theorem-unit-ball-compact-finite-dim}
A normed space $X$ is finite-dimensional if
and only if $B_X$ is compact.
\end{theorem}



\section{Hahn--Banach theorems}
\label{section-hahn-banach}

\begin{definition}
\label{definition-sublinear}
Let $E$ be a real vector space. A function
$p:E\to\mathbb{R}$ is {\it sublinear} if
$p(\lambda x)=\lambda p(x)$ for
$\lambda\geq0$ and $p(x+y)\leq p(x)+p(y)$ for
all $x,y \in E$.
\end{definition}

\begin{theorem}
[Hahn--Banach, real analytic form]
\label{theorem-hahn-banach-real}
Let $E$ be a real vector space, let
$p:E\to\mathbb{R}$ be sublinear, let
$G\subset E$ be a linear subspace, and let
$g:G\to\mathbb{R}$ be linear with $g(x)\leq
p(x)$ for all $x\in G$. Then there is a
linear functional $f:E\to\mathbb{R}$
extending $g$ such that $f(x)\leq p(x)$ for
all $x\in E$.
\end{theorem}


\begin{lemma}
\label{lemma-hahn-banach-one-dimensional-extension}
Under the hypotheses of the real
Hahn--Banach theorem, if
$x_0\in E\setminus G$, then $g$ extends to
$\widetilde g$ on
$\operatorname{span}(G\cup\{x_0\})$ with
$\widetilde g\leq p$.
\end{lemma}

\begin{proof}
Any such extension must have the form
$
\widetilde g(y+\alpha x_0)
=
g(y)+\alpha c.
$
For $\alpha=1$, the inequality
$\widetilde g\leq p$ requires
$
c\leq p(y+x_0)-g(y)
\qquad(y\in G).
$
Set
$
c=\inf_{y\in G}\{p(y+x_0)-g(y)\}.
$
This infimum is finite: for
$y_1,y_2\in G$,
\begin{align*}
g(y_1)-g(y_2)
&=g(y_1-y_2)\\
&\leq p(y_1-y_2)\\
&\leq p(y_1+x_0)+p(-y_2-x_0),
\end{align*}
so
$
-p(-y_2-x_0)-g(y_2)
\leq c
\leq p(y_1+x_0)-g(y_1).
$

For $\alpha>0$, use
$y_1=y/\alpha$ in the upper bound for $c$:
$
\widetilde g(y+\alpha x_0)
\leq p(y+\alpha x_0).
$
For $\alpha<0$, use the lower bound for $c$
and the same calculation with $|\alpha|$.
Thus $\widetilde g\leq p$ on the enlarged
subspace.
\end{proof}

\begin{proof}
[Proof of Theorem \ref{theorem-hahn-banach-real}]
Let $\mathbb P$ be the set of pairs $(F,h)$
such that
$
G\subseteq F\subseteq E,
$
$h:F\to\mathbb R$ is linear,
$h|_G=g$, and $h\leq p$ on $F$. Order
$\mathbb P$ by extension. A chain has an
upper bound obtained by taking the union of
the compatible functionals in the chain.
Hence Zorn's lemma gives a maximal element
$(F,h)$.

If $F\neq E$, choose $x_0\in E\setminus F$.
Lemma
\ref{lemma-hahn-banach-one-dimensional-extension}
extends $h$ to a strictly larger subspace,
contradicting maximality. Thus $F=E$.
\end{proof}

\begin{theorem}
[Hahn--Banach, complex analytic form]
\label{theorem-hahn-banach-complex}
Let $E$ be a complex vector space, let
$p:E\to\mathbb{R}$ be a seminorm, let
$G\subset E$ be a subspace, and let
$g:G\to\mathbb{C}$ be linear with $|g(x)|\leq
p(x)$ for all $x\in G$. Then there is a
linear functional $f:E\to\mathbb{C}$
extending $g$ such that $|f(x)|\leq p(x)$ for
all $x\in E$.
\end{theorem}

\begin{proof}
Let $g_1 = \operatorname{Re} g$. Then $g_1$
is a real linear functional on $G$ and
$g_1(x) \leq |g(x)| \leq p(x)$. By the real
Hahn-Banach theorem, there is a real linear
extension $f_1$ of $g_1$ to $E$ such that
$f_1(x) \leq p(x)$ for all $x \in E$. Define
$f(x) = f_1(x) - i f_1(ix)$. One checks that
$f$ is complex linear, extends $g$, and
satisfies $|f(x)| \leq p(x)$.
\end{proof}

\begin{proposition}
\label
{proposition-hahn-banach-norm-preserving-%
extension}
Let $X$ be a normed space, let $Y\subset X$
be a linear subspace, and let $f\in Y^*$.
Then there is $F\in X^*$ such that $F|_Y=f$
and $\|F\|=\|f\|$.
\end{proposition}

\begin{proposition}
\label{proposition-norming-functional}
Let $X$ be a normed space and let $x_0 \in X
\setminus \{0\}$. There exists $f \in X^*$
such that $\|f\|=1$ and $f(x_0) = \|x_0\|$.
\end{proposition}

\begin{proof}
Let $Y = \operatorname{span}\{x_0\}$ and
define $g(\lambda x_0) = \lambda \|x_0\|$.
Then $\|g\|=1$ on $Y$. Apply Proposition
\ref{proposition-hahn-banach-norm-preserving-%
extension}.
\end{proof}

\begin{proposition}
\label{proposition-dual-separates-points}
The dual space $X^*$ separates points of $X$:
if $x \neq y$, there exists $f \in X^*$ such
that $f(x) \neq f(y)$.
\end{proposition}

\begin{theorem}
[Hahn--Banach, first geometric form]
\label{theorem-hahn-banach-geometric-1}
Let $X$ be a normed real vector space and let
$A, B \subset X$ be nonempty disjoint convex
sets. If $A$ is open, there exist $f \in X^*
\setminus \{0\}$ and $\alpha \in \mathbb{R}$
such that $f(a) < \alpha \leq f(b)$ for all
$a \in A, b \in B$.
\end{theorem}

\begin{theorem}
[Hahn--Banach, second geometric form]
\label{theorem-hahn-banach-geometric-2}
Let $X$ be a normed real vector space and let
$A, B \subset X$ be nonempty disjoint convex
sets. If $A$ is closed and $B$ is compact,
there exist $f \in X^* \setminus \{0\}$ and
$\alpha_1, \alpha_2 \in \mathbb{R}$ such that
$$
f(a) \leq \alpha_1 < \alpha_2 \leq f(b)
$$
for all $a \in A, b \in B$.
\end{theorem}



\begin{proposition}
\label{proposition-hahn-banach-separate-point-subspace}
Let $X$ be a normed space, let
$Y\subseteq X$ be a closed linear subspace,
and let $x_0\notin Y$. Then there exists
$f\in X^*$ such that
$
f|_Y=0,\qquad
f(x_0)=d(x_0,Y),\qquad
\|f\|=1.
$
\end{proposition}

\begin{proof}
On
$F=\operatorname{span}(Y\cup\{x_0\})$ define
$
g(y+\alpha x_0)=
\alpha\,d(x_0,Y).
$
For $\alpha\neq0$,
\begin{align*}
|g(y+\alpha x_0)|
&=
|\alpha|d(x_0,Y)\\
&\leq
|\alpha|
\left\|x_0+\frac y\alpha\right\|\\
&=
\|y+\alpha x_0\|.
\end{align*}
Thus $\|g\|\leq1$. Choosing a sequence
$y_n\in Y$ with
$\|x_0-y_n\|\to d(x_0,Y)$ shows
$\|g\|\geq1$. Hahn--Banach extends $g$ to
the required $f$.
\end{proof}

\begin{theorem}
\label{theorem-dual-separable-implies-space-separable}
If $X$ is a Banach space and $X^*$ is
separable, then $X$ is separable.
\end{theorem}

\begin{proof}
Choose a dense sequence
$(f_n)\subset S_{X^*}$. For each $n$, choose
$x_n\in S_X$ such that
$
|f_n(x_n)|>1-\frac1n.
$
Put
$
Y=\overline{\operatorname{span}}\{x_n:n\in
\mathbb N\}.
$
If $Y\neq X$, choose $x\notin Y$. By
Proposition
\ref{proposition-hahn-banach-separate-point-subspace},
there is $f\in S_{X^*}$ with $f|_Y=0$.
Choose a subsequence $f_{n_k}\to f$ in norm.
Then
$
|f(x_{n_k})|
\geq
|f_{n_k}(x_{n_k})|
-\|f_{n_k}-f\|,
$
and the right-hand side tends to $1$.
But $x_{n_k}\in Y$ and $f|_Y=0$, a
contradiction. Thus $Y=X$, so $X$ is
separable.
\end{proof}

\section{Weak and weak star topologies}
\label{section-weak-topologies}

\begin{definition}
\label{definition-weak-topology}
Let $X$ be a normed space. The {\it weak
topology} on $X$, denoted $\sigma(X,X^*)$, is
the coarsest topology on $X$ such that every
$f \in X^*$ is continuous.
\end{definition}

\begin{definition}
\label{definition-weak-star-topology}
Let $X$ be a normed space. The {\it weak star
topology} on $X^*$, denoted $\sigma(X^*,X)$,
is the coarsest topology on $X^*$ such that
for every $x \in X$, the evaluation map
$\text{ev}_x : X^* \to \mathbb{K}$ given by
$\text{ev}_x(f) = f(x)$ is continuous.
\end{definition}

\begin{proposition}
\label{proposition-weak-vs-norm-topology}
The weak topology is coarser than the norm
topology. They coincide if and only if $\dim
X < \infty$.
\end{proposition}

\begin{theorem}[Banach--Alaoglu]
\label{theorem-banach-alaoglu}
The closed unit ball $B_{X^*}$ of the dual
space $X^*$ is compact in the weak star
topology $\sigma(X^*,X)$.
\end{theorem}

\begin{proof}
For each $x \in X$, let $D_x = \{ \lambda \in
\mathbb{K} : |\lambda| \leq \|x\| \}$. By
Tychonoff's theorem, the product space $D =
\prod_{x \in X} D_x$ is compact in the
product topology. We can view $B_{X^*}$ as a
subset of $D$ via the map $\Phi : f \mapsto
(f(x))_{x \in X}$. The topology on $D$
restricted to $\Phi(B_{X^*})$ is precisely
the weak star topology. One checks that
$\Phi(B_{X^*})$ is a closed subset of $D$ by
verifying that the limit of a net of linear
functionals is linear. Thus $B_{X^*}$ is
compact.
\end{proof}

\begin{theorem}[Goldstine]
\label{theorem-goldstine}
Let $X$ be a normed space and $J : X \to
X^{**}$ the canonical embedding. Then
$J(B_X)$ is weak star dense in $B_{X^{**}}$.
\end{theorem}

\begin{proof}
Suppose $J(B_X)$ is not weak star dense.
Since it is convex, by the Hahn-Banach
separation theorem (geometric form), there
exists a weak star continuous linear
functional on $X^{**}$ that separates a point
in $B_{X^{**}}$ from $J(B_X)$. Any weak star
continuous functional on $X^{**}$ is of the
form $\text{ev}_f$ for some $f \in X^*$. Thus
there exists $f \in X^*$ and $\alpha \in
\mathbb{R}$ such that $\operatorname{Re}
\xi(f) > \alpha$ for some $\xi \in
B_{X^{**}}$ and $\operatorname{Re} Jx(f) \leq
\alpha$ for all $x \in B_X$. But $\sup_{x \in
B_X} \operatorname{Re} f(x) = \|f\|$, so
$\alpha \geq \|f\|$. On the other hand,
$|\xi(f)| \leq \|\xi\| \|f\| \leq \|f\|$, so
$\alpha < \|f\|$, a contradiction.
\end{proof}

\begin{definition}
\label{definition-reflexive-space}
A Banach space $X$ is {\it reflexive} if the
canonical embedding $J : X \to X^{**}$ is
surjective.
\end{definition}

\begin{theorem}[Kakutani]
\label{theorem-kakutani-reflexivity}
A Banach space $X$ is reflexive if and only
if its closed unit ball $B_X$ is compact in
the weak topology.
\end{theorem}

\begin{theorem}[Eberlein--\v{S}mulyan]
\label{theorem-eberlein-smulyan}
Let $A$ be a subset of a Banach space $X$.
The following are equivalent:
\begin{enumerate}
\item $A$ is relatively weakly compact.
\item $A$ is relatively weakly sequentially
compact.
\end{enumerate}
\end{theorem}



\subsection{Reflexivity and weak-star continuous functionals}
\label{subsection-reflexivity-course}

\begin{proposition}
\label{proposition-canonical-embedding-isometry}
The canonical map
$
J:X\longrightarrow X^{**},
\qquad
(Jx)(f)=f(x),
$
is an isometry.
\end{proposition}

\begin{proof}
For $x\in X$,
$
\|Jx\|
=
\sup_{f\in B_{X^*}}|f(x)|
\leq\|x\|.
$
If $x\neq0$, Hahn--Banach gives
$f\in X^*$ with $\|f\|=1$ and
$f(x)=\|x\|$. Hence
$
\|Jx\|\geq|(Jx)(f)|=\|x\|.
$
\end{proof}

\begin{proposition}
\label{proposition-closed-subspace-reflexive}
If $X$ is reflexive and $Y\subseteq X$ is a
closed linear subspace, then $Y$ is
reflexive.
\end{proposition}

\begin{proof}
Let $\xi\in Y^{**}$. Define
$\bar\xi\in X^{**}$ by
$
\bar\xi(f)=\xi(f|_Y),
\qquad f\in X^*.
$
Since $X$ is reflexive, there is $x\in X$
such that $J_Xx=\bar\xi$.

We claim $x\in Y$. Otherwise Hahn--Banach
gives $f\in X^*$ with $f|_Y=0$ but
$f(x)\neq0$. Then
$
0=\xi(f|_Y)=\bar\xi(f)=(J_Xx)(f)=f(x),
$
a contradiction.

Now let $g\in Y^*$. Extend $g$ by
Hahn--Banach to $\bar g\in X^*$. Then
\begin{align*}
(J_Yx)(g)
&=g(x)=\bar g(x)\\
&=(J_Xx)(\bar g)
=\bar\xi(\bar g)
=\xi(g).
\end{align*}
Thus $J_Yx=\xi$.
\end{proof}

\begin{lemma}
\label{lemma-kernel-intersection-span-functionals}
Let $f,f_1,\ldots,f_n:X\to\mathbb R$ be
linear. If
$
\bigcap_{i=1}^n\ker f_i\subseteq\ker f,
$
then
$
f\in\operatorname{span}\{f_1,\ldots,f_n\}.
$
\end{lemma}

\begin{proof}
Define
$
L:X\longrightarrow\mathbb R^n,
\qquad
L(x)=(f_1(x),\ldots,f_n(x)).
$
The kernel hypothesis implies that $f$
factors through $L$: there is a linear
$g:\operatorname{Im}L\to\mathbb R$ such that
$f=g\circ L$. Extend $g$ linearly to
$\mathbb R^n$. Thus
$
g(t_1,\ldots,t_n)=
\sum_{i=1}^n\lambda_it_i
$
for some scalars $\lambda_i$, and therefore
$f=\sum_i\lambda_if_i$.
\end{proof}

\begin{proposition}
\label{proposition-weakstar-continuous-bidual-is-evaluation}
Let $\xi\in X^{**}$. If
$\xi:X^*\to\mathbb K$ is continuous for the
weak-star topology $\sigma(X^*,X)$, then
$\xi=Jx$ for some $x\in X$.
\end{proposition}

\begin{proof}
Continuity at zero gives
$x_1,\ldots,x_n\in X$ and $\varepsilon>0$
such that
$
\bigcap_{i=1}^n
\{f\in X^*:|f(x_i)|<\varepsilon\}
\subseteq
\xi^{-1}(B(0,1)).
$
By linearity, this implies
$
\bigcap_{i=1}^n\ker(Jx_i)\subseteq\ker\xi.
$
Lemma
\ref{lemma-kernel-intersection-span-functionals}
gives
$
\xi=\sum_{i=1}^n\lambda_iJx_i
=
J\left(\sum_{i=1}^n\lambda_ix_i\right).
$
\end{proof}

\begin{theorem}
\label{theorem-reflexive-iff-weak-weakstar-dual}
A normed space $X$ is reflexive if and only
if
$
\sigma(X^*,X)=\sigma(X^*,X^{**}).
$
\end{theorem}

\begin{proof}
If $X$ is reflexive, $J(X)=X^{**}$ and the
two topologies are the same.

Conversely, suppose the two topologies agree.
Every $\xi\in X^{**}$ is continuous for
$\sigma(X^*,X^{**})$, hence for
$\sigma(X^*,X)$. By Proposition
\ref{proposition-weakstar-continuous-bidual-is-evaluation},
$\xi\in J(X)$. Thus $J$ is onto.
\end{proof}

\subsection{Minkowski functionals and geometric Hahn--Banach}
\label{subsection-minkowski-geometric-hahn-banach}

\begin{definition}
\label{definition-topological-vector-space}
A {\it topological vector space} is a vector
space $X$ with a topology for which addition
and scalar multiplication are continuous and
points are closed.
\end{definition}

\begin{definition}
\label{definition-minkowski-functional}
Let $A\subseteq X$. Its {\it Minkowski
functional} is
$
\mu_A(x)
=
\inf\left\{t>0:\frac xt\in A\right\}.
$
A set $A$ is {\it absorbing} if for every
$x\in X$ there is $t>0$ with $x/t\in A$.
It is {\it balanced} if
$\lambda x\in A$ whenever $x\in A$ and
$|\lambda|\leq1$.
\end{definition}

\begin{proposition}
\label{proposition-minkowski-unit-ball}
For the unit ball
$A=B_X$ of a normed space,
$A$ is convex, balanced, and absorbing, and
$
\mu_A(x)=\|x\|.
$
\end{proposition}

\begin{proof}
If $t>\|x\|$, then $x/t\in A$, so
$\mu_A(x)\leq\|x\|$. If
$0<t<\|x\|$, then $x/t\notin A$, hence
$t\leq\mu_A(x)$. Therefore
$\mu_A(x)=\|x\|$.
\end{proof}

\begin{proposition}
\label{proposition-minkowski-convex-absorbing}
Let $A$ be convex and absorbing. Then
$\mu_A$ is sublinear and
$
\{x:\mu_A(x)<1\}\subseteq A
\subseteq\{x:\mu_A(x)\leq1\}.
$
If $A$ is balanced, then
$\mu_A(\lambda x)=|\lambda|\mu_A(x)$.
\end{proposition}

\begin{proof}
Let $t>\mu_A(x)$ and $s>\mu_A(y)$. Then
$x/t,y/s\in A$, and convexity gives
$
\frac{x+y}{t+s}
=
\frac{t}{t+s}\frac xt+
\frac{s}{t+s}\frac ys
\in A.
$
Hence
$\mu_A(x+y)\leq t+s$, and letting
$t\downarrow\mu_A(x)$ and
$s\downarrow\mu_A(y)$ proves
subadditivity. Positive homogeneity follows
directly from the definition, and balancedness
extends it to all scalars by absolute value.
The two inclusions follow from convexity and
the definition of the infimum.
\end{proof}

\begin{proof}
[Proof of Theorem
\ref{theorem-hahn-banach-geometric-1}]
Choose $a\in A$, $b\in B$, put
$x_0=a-b$, and define
$
C=-x_0+A-B.
$
Then $C$ is open, convex, absorbing, and
contains $0$, while $x_0\notin C$ after the
corresponding translation convention.
On $E=\operatorname{span}\{x_0\}$ define
$f(\alpha x_0)=\alpha$. The Minkowski
functional $\mu_C$ dominates $f$ on $E$.
By Hahn--Banach, $f$ extends to
$\widetilde f:X\to\mathbb R$ with
$\widetilde f\leq\mu_C$.

Since $\widetilde f$ is bounded on the open
set $C\cap(-C)$, it is continuous. Applying
the inequality to translated points of
$A-B$ gives
$
\widetilde f(a)<\widetilde f(b)
\qquad(a\in A,\ b\in B).
$
With
$
\alpha=\inf_{b\in B}\widetilde f(b),
$
one obtains
$
\widetilde f(a)<\alpha\leq\widetilde f(b).
$
\end{proof}

\begin{lemma}
\label{lemma-compact-closed-separated-neighborhoods}
Let $A,B$ be disjoint subsets of a
topological vector space, with $A$ compact
and $B$ closed. There is an open neighborhood
$V$ of $0$ such that
$
(A+V)\cap(B+V)=\varnothing.
$
\end{lemma}

\begin{proof}
For every $x\in A\subset B^c$, continuity of
addition gives a symmetric neighborhood
$V_x$ of $0$ such that
$
x+V_x+V_x+V_x\subseteq B^c.
$
Compactness gives
$x_1,\ldots,x_n\in A$ with
$
A\subseteq\bigcup_{i=1}^n(x_i+V_{x_i}).
$
Set $V=\bigcap_iV_{x_i}$. If a point belonged
to both $A+V$ and $B+V$, it would force some
$x_i+V_{x_i}+V_{x_i}+V_{x_i}$ to meet $B$,
contrary to the choice of $V_{x_i}$.
\end{proof}

\begin{remark}
\label{remark-geometric-hb-second-source}
The Portuguese course notes state the second
geometric Hahn--Banach theorem in the compact
versus closed case and prove the preceding
neighborhood-separation lemma; they leave the
final deduction of the theorem as an
exercise.
\end{remark}

\begin{corollary}
\label{corollary-convex-weak-norm-closure}
If $C$ is convex in a normed space $X$, then
$
\overline C^{\,\|\cdot\|}
=
\overline C^{\,w}.
$
\end{corollary}

\begin{proof}
Since the weak topology is coarser than the
norm topology,
$\overline C^{\,\|\cdot\|}
\subseteq\overline C^{\,w}$.
If
$x\notin\overline C^{\,\|\cdot\|}$, apply
geometric Hahn--Banach to the point
$\{x\}$ and the closed convex set
$\overline C^{\,\|\cdot\|}$. There is
$f\in X^*$ separating them. Hence a weak
open neighborhood of $x$ misses $C$, so
$x\notin\overline C^{\,w}$.
\end{proof}

\section
{Baire category and Banach--Steinhaus}
\label{section-baire-banach-steinhaus}

\begin{theorem}[Baire category theorem]
\label{theorem-baire}
Let $M$ be a non-empty complete metric space.
\begin{enumerate}
\item If $\{U_n\}_{n=1}^\infty$ is a sequence
of open dense subsets of $M$, then
$\bigcap_{n=1}^\infty U_n$ is dense in $M$.
\item $M$ is not a countable union of nowhere
dense sets.
\end{enumerate}
\end{theorem}

\begin{proof}
(1) Let $W \subseteq M$ be a non-empty open
set. Since $U_1$ is dense, there exists $x_1
\in W \cap U_1$ and $r_1 \in (0,1)$ such that
$\overline{B(x_1, r_1)} \subseteq W \cap
U_1$. Proceeding inductively, given $x_n,
r_n$, there exists $x_{n+1}, r_{n+1}$ such
that $\overline{B(x_{n+1}, r_{n+1})}
\subseteq B(x_n, r_n) \cap U_{n+1}$ and
$r_{n+1} < r_n / 2$. The sequence $(x_n)$ is
Cauchy, and its limit $x$ belongs to $\bigcap
\overline{B(x_n, r_n)}$, which is contained
in $W \cap (\bigcap U_n)$.
\end{proof}

\begin{theorem}
[Banach--Steinhaus / Uniform boundedness
principle]
\label{theorem-banach-steinhaus}
Let $X$ be a Banach space, $Y$ a normed
space, and let $\mathcal F \subset \mathcal
L(X,Y)$. If $\sup_{T \in \mathcal F} \|Tx\| <
\infty$ for each $x \in X$, then $\sup_{T \in
\mathcal F} \|T\| < \infty$.
\end{theorem}

\begin{proof}
For each $m \in \mathbb{N}$, let $X_m = \{ x
\in X : \|Tx\| \leq m \text{ for all } T \in
\mathcal F \}$. Each $X_m = \bigcap_{T \in
\mathcal F} \{x : \|Tx\| \leq m\}$ is closed.
By assumption, $X = \bigcup_{m=1}^\infty
X_m$. Since $X$ is Banach, by Baire's
theorem, some $X_{m_0}$ has non-empty
interior. Let $B(y, r) \subseteq X_{m_0}$.
For any $x \in X$ with $\|x\| \leq 1$, $y+rx
\in X_{m_0}$, so $\|T(y+rx)\| \leq m_0$. Then
$\|Trx\| \leq \|T(y+rx)\| + \|Ty\| \leq
2m_0$, so $\|Tx\| \leq 2m_0/r$. Thus $\|T\|
\leq 2m_0/r$ for all $T \in \mathcal F$.
\end{proof}

\begin{example}
\label{example-completeness-necessary-bs}
Let $X = c_{00}$ with the sup-norm. Define
$f_n \in X^*$ by $f_n(x) = n x_n$. For each
$x \in X$, $f_n(x) = 0$ for large $n$, so
$\sup_n |f_n(x)| < \infty$. However, $\|f_n\|
= n \to \infty$. Thus completeness of $X$ is
essential.
\end{example}

\begin{lemma}
\label{lemma-pointwise-limit-linear}
Let $X$ be a Banach space and $Y$ a normed
space. If $(T_n) \subset \mathcal L(X,Y)$ is
a sequence such that $Tx = \lim_n T_n x$
exists for all $x \in X$, then $T \in
\mathcal L(X,Y)$ and $\|T\| \leq \liminf_n
\|T_n\|$.
\end{lemma}

\section{Krein--Milman theorem}
\label{section-krein-milman}

\begin{definition}
\label{definition-extreme-point}
Let $K$ be a convex subset of a vector space.
\begin{enumerate}
\item A non-empty subset $L \subseteq K$ is
an {\it extreme set} of $K$ if for any $x, y
\in K$ and $\lambda \in (0,1)$, $\lambda x +
(1-\lambda)y \in L$ implies $x, y \in L$.
\item A point $x \in K$ is an {\it extreme
point} if $\{x\}$ is an extreme set. The set
of extreme points is denoted $E(K)$.
\end{enumerate}
\end{definition}

\begin{example}
\label{example-extreme-points}
\begin{enumerate}
\item $E(B_{c_0}) = \emptyset$.
\item $E(B_{\ell_\infty}) = \{ (x_n) : x_n
\in \{-1, 1\} \text{ for all } n \}$.
\item $E(B_{\ell_1}) = \{ \pm e_n : n \in
\mathbb{N} \}$.
\item For $1 < p < \infty$, $E(B_{\ell_p}) =
S_{\ell_p}$ (the unit sphere).
\end{enumerate}
\end{example}

\begin{theorem}[Krein--Milman]
\label{theorem-krein-milman}
Let $X$ be a locally convex topological
vector space and $K \subseteq X$ a non-empty
compact convex set. Then $K$ is the closed
convex hull of its extreme points: $K =
\overline{\operatorname{conv}}(E(K))$.
\end{theorem}

\begin{lemma}
\label{lemma-c0-not-dual}
Since $E(B_{c_0}) = \emptyset$, $c_0$ is not
isometrically isomorphic to the dual of any
Banach space (if it were, $B_{c_0}$ would be
weak* compact and would have extreme points
by Krein-Milman).
\end{lemma}


\begin{proof}
[Proof of Theorem \ref{theorem-krein-milman}]
The inclusion
$\overline{\operatorname{conv}}E(K)\subseteq
K$ is immediate because $K$ is closed and
convex.

Let $\mathbb P$ be the family of nonempty
extreme subsets of $K$, ordered by reverse
inclusion. For a fixed $L\in\mathbb P$,
Zorn's lemma gives a minimal extreme subset
$S\subseteq L$: chains have lower bounds
given by intersection, which is nonempty by
compactness.

We claim $S$ is a singleton. If it contained
two distinct points, geometric Hahn--Banach
would give $f\in X^*$ nonconstant on $S$.
Since $S$ is compact,
$
\mu=\max_{x\in S}f(x)
$
exists, and
$
S_f=\{x\in S:f(x)=\mu\}
$
is a nonempty proper subset of $S$. Moreover
$S_f$ is extreme in $K$: if
$\lambda x+(1-\lambda)y\in S_f$ with
$x,y\in K$, extremality of $S$ first gives
$x,y\in S$, and maximality of $f$ on $S$
then forces $f(x)=f(y)=\mu$. This contradicts
minimality of $S$. Hence every nonempty
extreme subset contains an extreme point of
$K$.

Suppose now that
$x_0\in K\setminus
\overline{\operatorname{conv}}E(K)$.
Geometric Hahn--Banach gives
$f\in X^*$ such that
$
\sup_{\overline{\operatorname{conv}}E(K)}f
<
f(x_0).
$
The set
$
K_f=
\left\{x\in K:
f(x)=\max_{z\in K}f(z)\right\}
$
is a nonempty extreme subset of $K$, hence
contains some $y\in E(K)$. But then
$
f(y)
\leq
\sup_{E(K)}f
<
f(x_0)
\leq
f(y),
$
a contradiction.
\end{proof}

\begin{corollary}
\label{corollary-dual-ball-extreme-point}
If $X$ is a Banach space, then
$E(B_{X^*})\neq\varnothing$.
\end{corollary}

\section
{Open mapping and closed graph theorems}
\label{section-open-mapping-closed-graph}

\begin{theorem}[Open mapping theorem]
\label{theorem-open-mapping}
Let $X, Y$ be Banach spaces and let $T \in
\mathcal L(X,Y)$ be surjective. Then $T$ is
an open map.
\end{theorem}

\begin{proof}
It is enough to show that
$
0\in\operatorname{int}T(rB_X)
\qquad\text{for every }r>0.
$
Since
$
Y=T(X)=\bigcup_{n=1}^\infty T(nB_X),
$
Baire's theorem gives
$
\operatorname{int}\overline{T(rB_X)}
\neq\varnothing
$
for every $r>0$ after rescaling.

Next,
$
\overline{T(rB_X/2)}
-
\overline{T(rB_X/2)}
\subseteq
\overline{T(rB_X)}.
$
If $U$ is a nonempty open subset of
$\overline{T(rB_X/2)}$, then
$0\in U-U$, hence
$
0\in\operatorname{int}\overline{T(rB_X)}.
$

Finally we prove
$
\overline{T(rB_X/2)}
\subseteq T(rB_X).
$
Fix
$y_1=y\in\overline{T(rB_X/2)}$. Construct
$x_n\in (r/2^n)B_X$ and
$y_n\in\overline{T((r/2^n)B_X)}$ so that
$
y_{n+1}=y_n-Tx_n.
$
This is possible because
$0$ is an interior point of the relevant
closures. Since
$\sum_n\|x_n\|<r$, the series
$x=\sum_nx_n$ converges in $X$ and belongs
to $rB_X$. Moreover
$
\sum_{n=1}^NTx_n
=
y_1-y_{N+1}
\longrightarrow y,
$
because $y_n\to0$. Thus $Tx=y$.

Therefore $0$ is an interior point of
$T(rB_X)$ for every $r>0$. If $U\subset X$
is open and $x\in U$, choose
$x+rB_X\subset U$. A neighborhood of zero in
$T(rB_X)$ translates to a neighborhood of
$Tx$ contained in $T(U)$, so $T(U)$ is open.
\end{proof}

\begin{proposition}[Bounded inverse theorem]
\label{proposition-bounded-inverse}
If $T \in \mathcal L(X,Y)$ is a bijection
between Banach spaces, then $T^{-1}$ is
bounded.
\end{proposition}

\begin{theorem}[Closed graph theorem]
\label{theorem-closed-graph}
Let $X, Y$ be Banach spaces and let $T : X
\to Y$ be a linear map. If the graph
$\Gamma(T) = \{ (x, Tx) : x \in X \}$ is
closed in $X \times Y$, then $T$ is bounded.
\end{theorem}

\begin{proof}
Consider the projection $\pi_1 : \Gamma(T)
\to X$ given by $(x, Tx) \mapsto x$.
$\Gamma(T)$ is a Banach space as it is a
closed subspace of $X \times Y$. $\pi_1$ is a
bounded linear bijection. By the Bounded
Inverse Theorem, $\pi_1^{-1}$ is bounded.
Since $T = \pi_2 \circ \pi_1^{-1}$ (where
$\pi_2$ is the projection on $Y$), $T$ is
bounded.
\end{proof}

\section{Complemented subspaces}
\label{section-complemented-subspaces}

\begin{definition}
\label{definition-complemented-subspace}
A closed subspace $Y \subseteq X$ is {\it
complemented} if there exists a closed
subspace $Z \subseteq X$ such that $X = Y
\oplus Z$.
\end{definition}

\begin{proposition}
\label{proposition-finite-codim-complemented}
If $Y$ is a closed subspace of a Banach space
$X$ and $\operatorname{codim} Y < \infty$,
then $Y$ is complemented.
\end{proposition}

\begin{theorem}[Phillips]
\label{theorem-c0-not-complemented}
The subspace $c_0$ is not complemented in
$\ell_\infty$.
\end{theorem}

\begin{definition}
\label{definition-almost-disjoint-family}
A family $(A_i)_{i\in I}$ of subsets of
$\mathbb N$ is {\it almost disjoint} if
$
|A_i\cap A_j|<\infty
\qquad(i\neq j).
$
\end{definition}

\begin{proposition}
\label{proposition-uncountable-almost-disjoint-family}
There exists an almost-disjoint family of
infinite subsets of $\mathbb N$ of
cardinality $2^{\aleph_0}$.
\end{proposition}

\begin{proof}
For every irrational $r$, choose a sequence
of distinct rationals converging to $r$.
After fixing a bijection
$\mathbb Q\to\mathbb N$, take the image of
each such sequence. Two of the resulting
sets can have only finitely many common
terms, since their corresponding rational
sequences converge to distinct limits.
\end{proof}

\begin{proposition}
\label{proposition-operator-vanishing-c0-vanishes-large-linfty}
Let $T:\ell_\infty\to\ell_\infty$ be bounded
and suppose $T|_{c_0}=0$. Then there is an
infinite $A\subseteq\mathbb N$ such that
$
T|_{\ell_\infty(A)}=0,
$
where $\ell_\infty(A)$ is extended by zero
outside $A$.
\end{proposition}

\begin{proof}
Let $(A_i)_{i\in I}$ be an almost-disjoint
family as above. If the conclusion were
false, choose
$x_i\in S_{\ell_\infty(A_i)}$ with
$Tx_i\neq0$. Since $|I|=2^{\aleph_0}$,
there are a coordinate $n$, a number
$\delta>0$, and an uncountable
$J\subseteq I$ such that, after replacing
$x_i$ by $-x_i$ when necessary,
$
e_n^*(Tx_i)>\delta
\qquad(i\in J).
$
For finite $F\subseteq J$,
$
\left\|
T\left(\sum_{i\in F}x_i\right)
\right\|
\geq |F|\delta.
$
Almost-disjointness lets us write
$
\sum_{i\in F}x_i=u+v,
$
where $u$ has finite support and
$\|v\|_\infty\leq1$. Since $Tu=0$,
$
|F|\delta
\leq\|Tv\|
\leq\|T\|,
$
a contradiction for large finite $F$.
\end{proof}

\begin{proof}[Proof of Theorem
\ref{theorem-c0-not-complemented}]
If a bounded projection
$P:\ell_\infty\to c_0$ existed, then
$
(I-P)|_{c_0}=0.
$
By Proposition
\ref{proposition-operator-vanishing-c0-vanishes-large-linfty},
there would be an infinite
$A\subseteq\mathbb N$ such that
$(I-P)|_{\ell_\infty(A)}=0$. Thus
$P$ would be the identity on
$\ell_\infty(A)$, forcing
$\ell_\infty(A)\subseteq c_0$, a
contradiction.
\end{proof}


\begin{definition}
\label{definition-compact-operator}
Let $X, Y$ be normed spaces. $T \in \mathcal
L(X,Y)$ is {\it compact} if $T(B_X)$ has
compact closure in $Y$. The set of compact
operators is denoted $\mathcal K(X,Y)$.
\end{definition}

\begin{proposition}
\label{proposition-compact-operator-ideal}
$\mathcal K(X)$ is a closed two-sided ideal
in $\mathcal L(X)$.
\end{proposition}

\begin{theorem}[Schauder]
\label{theorem-schauder-compact-adjoint}
An operator $T \in \mathcal L(X,Y)$ is
compact if and only if its adjoint $T^* \in
\mathcal L(Y^*, X^*)$ is compact.
\end{theorem}

\begin{definition}
\label{definition-fredholm-operator}
$T \in \mathcal L(X,Y)$ is a {\it Fredholm
operator} if $\dim \ker T < \infty$ and
$\operatorname{codim} \operatorname{Im} T <
\infty$. The {\it index} of $T$ is
$$
\operatorname{ind}(T) = \dim \ker T -
\operatorname{codim} \operatorname{Im} T.
$$
\end{definition}

\begin{theorem}[Atkinson]
\label{theorem-atkinson}
Let $X$ be an infinite-dimensional Banach
space. $T \in \mathcal L(X)$ is Fredholm if
and only if its image in the Calkin algebra
$\mathcal L(X) / \mathcal K(X)$ is
invertible.
\end{theorem}

\begin{definition}
\label{definition-adjoint-banach}
For $T \in \mathcal L(X,Y)$, the {\it
adjoint} $T^* : Y^* \to X^*$ is defined by
$(T^*f)(x) = f(Tx)$.
\end{definition}

\begin{proposition}
\label{proposition-adjoint-properties}
The map $T \mapsto T^*$ is a linear isometry
from $\mathcal L(X,Y)$ to $\mathcal L(Y^*,
X^*)$. Moreover, $(ST)^* = T^*S^*$.
\end{proposition}



\begin{definition}
\label{definition-projection-banach}
A bounded operator $P:X\to X$ is a
{\it projection} if $P^2=P$.
\end{definition}

\begin{proposition}
\label{proposition-complemented-iff-projection}
A closed subspace $Y$ of a Banach space $X$
is complemented if and only if
$
Y=\operatorname{Im}P
$
for some bounded projection $P:X\to X$.
\end{proposition}

\begin{proof}
If $P$ is a projection, then
$
Y=\operatorname{Im}P=\ker(I-P),
\qquad
Z=\operatorname{Im}(I-P)=\ker P
$
are closed and $X=Y\oplus Z$.

Conversely, let $X=Y\oplus Z$ with $Y,Z$
closed. Algebraically define
$P(y+z)=y$. Let
$
\pi:X\to X/Z
$
be the quotient map. The restriction
$\pi|_Y:Y\to X/Z$ is a bounded bijection
between Banach spaces, so the bounded inverse
theorem gives continuity of its inverse.
Since
$
P=(\pi|_Y)^{-1}\circ\pi,
$
$P$ is bounded.
\end{proof}

\begin{proposition}
\label{proposition-finite-dimensional-complemented}
Every finite-dimensional subspace of a
Banach space is complemented.
\end{proposition}

\begin{proof}
Let $y_1,\ldots,y_n$ be a basis of $Y$, and
let $\widetilde y_i^*\in Y^*$ be the
coordinate functionals. Extend them by
Hahn--Banach to $y_i^*\in X^*$ and put
$
Z=\bigcap_{i=1}^n\ker y_i^*.
$
Then $Z$ is closed and
$
x=
\sum_{i=1}^ny_i^*(x)y_i+
\left(
x-\sum_{i=1}^ny_i^*(x)y_i
\right)
$
gives $X=Y\oplus Z$.
\end{proof}

\begin{proposition}
\label{proposition-weakstar-continuous-operators-are-adjoints}
Let $X,Y$ be Banach spaces and let
$S:Y^*\to X^*$ be linear. Then $S$ is
weak-star continuous if and only if there is
$T\in\mathcal L(X,Y)$ such that
$
S=T^*.
$
\end{proposition}

\begin{proof}
If $S=T^*$, weak-star continuity follows
from
$
(T^*f)(x)=f(Tx).
$

Conversely, suppose $S$ is weak-star
continuous. For each $x\in X$ define
$
\xi_x:Y^*\to\mathbb K,
\qquad
\xi_x(f)=(Sf)(x).
$
Then $\xi_x$ is weak-star continuous, so by
Proposition
\ref{proposition-weakstar-continuous-bidual-is-evaluation}
there is a unique $y\in Y$ with
$\xi_x=Jy$. Define $Tx=y$. The injectivity
of $J_Y$ shows $T$ is linear. Moreover
$
\|Tx\|=\|\xi_x\|\leq\|S\|\,\|x\|,
$
so $T$ is bounded. Finally,
$
(T^*f)(x)=f(Tx)
=(J_YTx)(f)
=\xi_x(f)
=(Sf)(x),
$
hence $S=T^*$.
\end{proof}

\section{Universal Banach spaces}
\label{section-universal-banach-spaces}

\begin{theorem}
\label{theorem-separable-embeds-linfty}
Every separable normed space $X$ is
isometrically isomorphic to a subspace of
$\ell_\infty$.
\end{theorem}

\begin{proof}
Let $\{x_n\}_{n=1}^\infty$ be a dense subset
of the unit sphere $S_X$. By Hahn-Banach,
there exist $f_n \in S_{X^*}$ such that
$f_n(x_n) = 1$. Define $T : X \to
\ell_\infty$ by $Tx = (f_n(x))_{n=1}^\infty$.
Since $|f_n(x)| \leq \|x\|$, $T$ is
well-defined and $\|Tx\| \leq \|x\|$. For any
$\varepsilon > 0$, choose $n$ such that $\|x
- \|x\|x_n\| < \varepsilon$. Then $\|x\| =
f_n(\|x\|x_n) < f_n(x) + \varepsilon \leq
\|Tx\| + \varepsilon$. Thus $\|Tx\| = \|x\|$.
\end{proof}

\begin{theorem}[Banach--Mazur]
\label{theorem-banach-mazur}
Every separable Banach space is isometrically
isomorphic to a subspace of $C[0,1]$.
\end{theorem}

\begin{proof}
The course notes use the Cantor space
$
\Delta=\{0,1\}^{\mathbb N}.
$
By Banach--Alaoglu,
$B_{X^*}$ is weak-star compact. Since $X$ is
separable, this compact space is metrizable.
For $x\in X$, define
$
(\varphi x)(f)=f(x),
\qquad f\in B_{X^*}.
$
Then
$
\varphi:X\to C(B_{X^*})
$
is a linear isometry.

Every compact metric space is a continuous
image of $\Delta$, so choose a continuous
surjection
$
\alpha:\Delta\to B_{X^*}.
$
Composition with $\alpha$ gives an isometry
$
\psi:C(B_{X^*})\to C(\Delta),
\qquad
(\psi h)(t)=h(\alpha(t)).
$
Thus $\psi\circ\varphi$ embeds $X$
isometrically in $C(\Delta)$. The usual
identification of $C(\Delta)$ with a
subspace of $C[0,1]$ yields the stated form.
\end{proof}

\begin{theorem}
[Pe\l{}czy\'{n}ski decomposition]
\label{theorem-pelczynski-decomposition}
Let $X, Y$ be Banach spaces such that each is
isomorphic to a complemented subspace of the
other.
\begin{enumerate}
\item If $X \cong X \oplus X$ and $Y \cong Y
\oplus Y$, then $X \cong Y$.
\item If $X \cong \ell_p(X)$ for some $1 \leq
p \leq \infty$, then $X \cong Y$.
\end{enumerate}
\end{theorem}

\begin{proof}
Choose closed spaces $E,F$ such that
$
X\cong Y\oplus E,
\qquad
Y\cong X\oplus F.
$
For (1), using
$X\cong X\oplus X$ and
$Y\cong Y\oplus Y$,
\begin{align*}
X
&\cong Y\oplus E\\
&\cong Y\oplus Y\oplus E\\
&\cong X\oplus F\oplus Y\oplus E\\
&\cong X\oplus F\oplus X\\
&\cong X\oplus F\\
&\cong Y.
\end{align*}

For (2), suppose $X\cong\ell_p(X)$. First,
\begin{align*}
Y
&\cong X\oplus F\\
&\cong X\oplus\ell_p(X)\oplus F\\
&\cong Y\oplus X.
\end{align*}
Then
\begin{align*}
X
&\cong\ell_p(X)\\
&\cong\ell_p(Y\oplus E)\\
&\cong\ell_p(Y)\oplus\ell_p(E)\\
&\cong Y\oplus\ell_p(Y)\oplus\ell_p(E)\\
&\cong Y.
\end{align*}
The same argument is recorded in the source
for $c_0(X)$.
\end{proof>



\section{Representation of functionals on $C(K)$}
\label{section-riesz-markov-course}

Let $X$ be a compact metric space and let
$\mathcal B_X$ be its Borel $\sigma$-algebra.

\begin{definition}
\label{definition-signed-borel-measure-total-variation}
A finite Borel measure allowed to take
positive and negative values is called a
{\it signed measure}. Its {\it total
variation} is
$
|\mu|(B)
=
\sup
\left\{
\sum_{n=1}^\infty|\mu(B_n)|:
B=\bigsqcup_{n=1}^\infty B_n,\ 
B_n\in\mathcal B_X
\right\}.
$
Write
$
M(X)
=
\{\mu:\mu\text{ is a finite signed Borel
measure on }X\}
$
and
$
\|\mu\|_M=|\mu|(X).
$
\end{definition}

\begin{definition}
\label{definition-positive-functional-cx}
A functional $\varphi\in C(X)^*$ is
{\it positive} if
$
f\geq0\quad\Longrightarrow\quad
\varphi(f)\geq0.
$
\end{definition}

\begin{remark}
\label{remark-positive-functional-norm}
If $\varphi$ is positive, then
$
\|\varphi\|=\varphi(1).
$
Indeed, for $f\in C(X)$,
$
\|f\|_\infty\pm f\geq0,
$
so
$
|\varphi(f)|
\leq
\varphi(1)\|f\|_\infty.
$
Equality in norm follows by evaluating at
the constant function $1$.
\end{remark}

\begin{theorem}
\label{theorem-positive-riesz-representation-course}
Let $X$ be a compact metric space and let
$\varphi\in C(X)^*$ be positive. Then there
is a unique finite positive Borel measure
$\mu$ on $X$ such that
$
\varphi(f)=\int_X f\,d\mu
\qquad\text{for all }f\in C(X).
$
\end{theorem}

\begin{proof}
For an open set $A\subseteq X$, define
$
r_\varphi(A)
=
\sup
\left\{
\varphi(f):
\operatorname{supp}f\subseteq A,\ 
0\leq f\leq1
\right\}.
$
For an arbitrary $B\subseteq X$, define
$
\mu^*(B)
=
\inf
\{r_\varphi(A):B\subseteq A,\ A\text{
open}\}.
$
The course notes construct the representing
measure from this outer measure.

Clearly $\mu^*(\varnothing)=0$, and
monotonicity follows directly from the
definition. To prove countable subadditivity,
first let $B=\bigcup_iB_i$ be a union of open
sets and let $f\in C(X)$ satisfy
$\operatorname{supp}f\subseteq B$ and
$0\leq f\leq1$. Compactness of
$\operatorname{supp}f$ gives a finite
subcover $B_1,\ldots,B_N$. Choose a
partition of unity
$\eta_1,\ldots,\eta_N$ subordinate to this
finite cover. Then
$
f=\sum_{i=1}^N f\eta_i,
$
with
$\operatorname{supp}(f\eta_i)\subseteq B_i$
and $0\leq f\eta_i\leq1$, so
$
\varphi(f)
\leq
\sum_{i=1}^Nr_\varphi(B_i).
$
Taking the supremum over $f$ gives
$
r_\varphi(B)
\leq
\sum_i r_\varphi(B_i).
$

For arbitrary sets $E_i$, choose open
$B_i\supseteq E_i$ with
$
r_\varphi(B_i)
\leq
\mu^*(E_i)+\varepsilon2^{-i}.
$
Then
$
\mu^*\left(\bigcup_iE_i\right)
\leq
\sum_i\mu^*(E_i)+\varepsilon.
$
Letting $\varepsilon\downarrow0$ proves
countable subadditivity.

The notes next verify that $\mu^*$ is a
metric outer measure: if
$d(E_1,E_2)>0$, choose disjoint open
neighborhoods of $E_1,E_2$ and obtain
$
\mu^*(E_1\cup E_2)
=
\mu^*(E_1)+\mu^*(E_2).
$
Therefore the restriction of $\mu^*$ to the
Borel sets is a finite Borel measure $\mu$.
Moreover,
$
\mu(X)=\mu^*(X)=\varphi(1).
$

It remains to show that $\mu$ represents
$\varphi$. By linearity it is enough to
consider $0\leq f\leq1$. The notes divide the
range of $f$ into $N$ intervals and construct
continuous functions $f_1,\ldots,f_N$ such
that
$
f=\sum_{n=1}^N f_n,
$
each $Nf_n$ is between $0$ and $1$, equals
$1$ on a slightly smaller superlevel set,
and is supported in a slightly larger
superlevel set. From the definition of
$r_\varphi$ one obtains upper and lower
Riemann-sum bounds for both
$\varphi(f)$ and $\int f\,d\mu$. The
difference is bounded by
$
\frac{2\mu(X)}N.
$
Letting $N\to\infty$ gives
$
\varphi(f)=\int_Xf\,d\mu.
$

For uniqueness, suppose $\mu'$ is another
positive finite Borel measure representing
$\varphi$. If $B$ is open and
$0\leq f\leq1$ has
$\operatorname{supp}f\subseteq B$, then
$
\varphi(f)=\int_Bf\,d\mu'\leq\mu'(B).
$
Taking the supremum over such $f$ gives
$\mu(B)\leq\mu'(B)$. Conversely, regularity
of $\mu'$ lets one choose compact
$K\subseteq B$ with
$\mu'(B)\leq\mu'(K)+\varepsilon$ and a
continuous $f$ with
$f|_K=1$, $0\leq f\leq1$, and
$\operatorname{supp}f\subseteq B$. Then
$
\mu'(B)
\leq
\varphi(f)+\varepsilon
\leq
\mu(B)+\varepsilon.
$
Thus the measures agree on open sets, hence
on all Borel sets.
\end{proof}

\begin{proposition}
\label{proposition-functional-positive-negative-parts}
Let $\varphi\in C(X)^*$. There are positive
functionals $\varphi^+,\varphi^-\in C(X)^*$
such that
$
\varphi=\varphi^+-\varphi^-,
\qquad
\|\varphi\|
=
\varphi^+(1)+\varphi^-(1).
$
\end{proposition}

\begin{theorem}[Riesz representation]
\label{theorem-riesz-representation-ck}
Let $X$ be a compact metric space. For every
$\varphi\in C(X)^*$ there is a unique
$\mu\in M(X)$ such that
$
\varphi(f)=\int_Xf\,d\mu
\qquad
\text{for all }f\in C(X),
$
and
$
\|\varphi\|=|\mu|(X)=\|\mu\|_M.
$
Thus
$
C(X)^*\cong M(X)
$
isometrically.
\end{theorem}

\begin{proof}
Write
$
\varphi=\varphi^+-\varphi^-
$
as in Proposition
\ref{proposition-functional-positive-negative-parts}.
By Theorem
\ref{theorem-positive-riesz-representation-course}
there are positive measures
$\mu_+,\mu_-$ representing
$\varphi^+,\varphi^-$. Put
$
\mu=\mu_+-\mu_-.
$
Then $\mu$ represents $\varphi$, and
\begin{align*}
|\varphi(f)|
&=
\left|\int_Xf\,d\mu\right|\\
&\leq
\int_X|f|\,d|\mu|\\
&\leq
\|f\|_\infty|\mu|(X),
\end{align*}
so
$\|\varphi\|\leq|\mu|(X)$. The decomposition
also gives
$
|\mu|(X)
\leq
\mu_+(X)+\mu_-(X)
=
\varphi^+(1)+\varphi^-(1)
=
\|\varphi\|,
$
hence equality.

If $\mu'$ also represents $\varphi$, then
$\nu=\mu-\mu'$ integrates every continuous
function to zero. Writing
$
\nu=\nu^+-\nu^-,
\qquad
\nu^\pm=\frac12(|\nu|\pm\nu),
$
the positive measures $\nu^+$ and $\nu^-$
define the same positive functional.
Uniqueness in the positive case gives
$\nu^+=\nu^-$, hence $\nu=0$.
\end{proof}


\section{Lyapunov convexity and fixed points}
\label{section-lyapunov-convexity}

\begin{definition}
\label{definition-nonatomic-signed-measure}
A signed measure $\mu$ is {\it nonatomic} if
whenever $|\mu|(A)>0$, there exists a
measurable $B\subseteq A$ such that
$
0<|\mu|(B)<|\mu|(A).
$
\end{definition}

\begin{theorem}[Lyapunov convexity theorem]
\label{theorem-lyapunov-convexity}
Let $\mu_1,\ldots,\mu_n$ be finite nonatomic
signed measures on the same measurable
space. Then the range of the vector measure
$
\mu(A)=
(\mu_1(A),\ldots,\mu_n(A))
$
is a compact convex subset of $\mathbb R^n$.
\end{theorem}

\begin{proof}
Put
$
\nu=|\mu_1|+\cdots+|\mu_n|.
$
Define
$
\Lambda:L^\infty(\nu)\longrightarrow
\mathbb R^n,
\qquad
\Lambda(f)=
\left(
\int f\,d\mu_1,\ldots,
\int f\,d\mu_n
\right).
$
By Radon--Nikodym, for each $i$ there is
$g_i\in L^1(\nu)$ such that
$
\int f\,d\mu_i
=
\int fg_i\,d\nu.
$
Hence $\Lambda$ is weak-star continuous on
$L^\infty(\nu)=(L^1(\nu))^*$.

Consider
$
K=
\{f\in L^\infty(\nu):0\leq f\leq1\}.
$
The notes record that $K$ is weak-star
closed. By Banach--Alaoglu it is weak-star
compact, and it is convex. Therefore
$\Lambda(K)$ is compact and convex.

Every value of the vector measure belongs to
$\Lambda(K)$ because
$
\mu(A)=\Lambda(\chi_A).
$
Conversely, fix $\xi\in\Lambda(K)$ and set
$
K_\xi=
\Lambda^{-1}(\{\xi\})\cap K.
$
This is nonempty, convex, and weak-star
compact. By Krein--Milman it has an extreme
point $f$.

We claim that $f$ is a characteristic
function. Otherwise there are a measurable
set $A$ of positive $\nu$-measure and
$r>0$ such that
$
r\leq f\leq1-r
\qquad\text{on }A.
$
Because $\nu$ is nonatomic,
$L^\infty(A,\nu)$ is infinite-dimensional,
whereas $\ker\Lambda$ has finite codimension.
Thus there is a nonzero
$
g\in L^\infty(A,\nu)\cap\ker\Lambda
$
with $\|g\|_\infty<r$. Then
$f\pm g\in K_\xi$ and
$
f=\frac12(f+g)+\frac12(f-g),
$
contradicting extremality. Hence
$f=\chi_A$ for some measurable $A$, and
$\xi=\mu(A)$.
\end{proof}

\begin{theorem}[Markov--Kakutani]
\label{theorem-markov-kakutani}
Let $X$ be a Banach space and
$K\subseteq X$ a nonempty convex weakly
compact set. Let $\mathcal T$ be a commuting
family of weakly continuous affine maps
$T:K\to K$. Then there exists $x\in K$ such
that
$
Tx=x
\qquad\text{for every }T\in\mathcal T.
$
\end{theorem}

\begin{proof}
First consider one map $T$. Define
$
A=\{(x,x):x\in K\},
\qquad
B=\{(x,Tx):x\in K\}
$
inside $X\times X$. If $T$ had no fixed
point, $A$ and $B$ would be disjoint compact
convex sets. Geometric Hahn--Banach would
give $F\in(X\times X)^*$ and
$\alpha<\beta$ such that
$
F(x,x)<\alpha<\beta<F(x,Tx)
\qquad(x\in K).
$
Hence
$
F(0,Tx)-F(0,x)>\beta-\alpha.
$
Applying this to
$x,T x,\ldots,T^{n-1}x$ and summing
telescopically gives
$
F(0,T^nx)-F(0,x)>n(\beta-\alpha),
$
contradicting boundedness of the continuous
functional $F$ on the compact set
$\{0\}\times K$. Thus every $T$ has a fixed
point.

For each $T\in\mathcal T$, let
$
K_T=\{x\in K:Tx=x\}.
$
Each $K_T$ is nonempty compact and convex.
For a finite family
$T_1,\ldots,T_n$, commutativity implies
$T_j$ preserves the common fixed-point set
of $T_1,\ldots,T_{j-1}$. Applying the
one-map case inductively shows
$
K_{T_1}\cap\cdots\cap K_{T_n}\neq\varnothing.
$
Compactness and the finite intersection
property give
$
\bigcap_{T\in\mathcal T}K_T\neq\varnothing.
$
\end{proof}


\section{Amenable groups}
\label{section-amenable-groups}

\begin{definition}
\label{definition-mean-set}
A {\it mean} on a set $X$ is a positive
functional
$
m\in\ell_\infty(X)^*
$
with $\|m\|=1$.
\end{definition}

\begin{remark}
\label{remark-mean-finitely-additive-probability}
A mean can be viewed as a finitely additive
probability measure on $X$.
\end{remark}

\begin{definition}
\label{definition-amenable-group}
Let $G$ be a group. It acts on
$\ell_\infty(G)$ by
$
(g\cdot a)(h)=a(g^{-1}h).
$
A mean $m$ is {\it invariant} if
$
m(g\cdot a)=m(a)
\qquad
\text{for all }a\in\ell_\infty(G),\ g\in G.
$
The group $G$ is {\it amenable} if it admits
an invariant mean.
\end{definition}

\begin{example}
\label{example-amenable-groups}
\begin{enumerate}
\item Finite groups are amenable, using
normalized counting measure.
\item The free group $\mathbb F_n$ is not
amenable for $n>1$.
\item The course notes record
$\mathbb Z$ as amenable.
\end{enumerate}
\end{example}

\begin{proof}
[Why $\mathbb F_2$ is not amenable]
Write
$
\mathbb F_2=\langle a,b\rangle
$
and, for
$c\in\{a,b,a^{-1},b^{-1}\}$, let
$\mathbb F_2^c$ be the set of reduced words
beginning with $c$. Then
$
a^{-1}\mathbb F_2^a
=
\mathbb F_2^a
\sqcup
\mathbb F_2^b
\sqcup
\mathbb F_2^{b^{-1}}
\sqcup\{e\}.
$
If an invariant mean $m$ existed, invariance
and finite additivity would force
$
m(\mathbb F_2^b
\sqcup\mathbb F_2^{b^{-1}}\sqcup\{e\})=0.
$
The analogous identity with $b$ forces the
remaining two first-letter pieces to have
total mean zero. Hence
$m(\mathbb F_2)=0$, contradicting
$m(1)=1$.
\end{proof}

\begin{theorem}
\label{theorem-abelian-groups-amenable}
Every abelian group is amenable.
\end{theorem}

\begin{proof}
Let
$
M(G)\subseteq\ell_\infty(G)^*
$
be the set of means. The course notes leave
as an exercise the verification that $M(G)$
is convex and weak-star compact.

For $g\in G$, define
$
T_g:M(G)\to M(G),
\qquad
(T_gm)(a)=m(g\cdot a).
$
Each $T_g$ is affine and weak-star
continuous. Since $G$ is abelian, the maps
$T_g$ commute. Markov--Kakutani gives a mean
$m\in M(G)$ with
$
T_gm=m
\qquad\text{for every }g\in G.
$
This is exactly invariance of the mean.
\end{proof}


\section{Hilbert spaces}
\label{section-hilbert-spaces}

\begin{definition}
\label{definition-inner-product-space}
An {\it inner product space} over
$\mathbb{K}$ is a vector space $H$ with a
function $\langle\cdot,\cdot\rangle:H\times
H\to\mathbb{K}$ which is linear in the first
variable, conjugate-linear in the second
variable, positive definite, and Hermitian:
$$
\langle x,y\rangle=\overline{\langle
y,x\rangle}.
$$
The associated norm is $\|x\|=\langle
x,x\rangle^{1/2}$.
\end{definition}

\begin{proposition}[Pythagoras]
\label{proposition-pythagoras}
If $x\perp y$, then
$$
\|x+y\|^2=\|x\|^2+\|y\|^2.
$$
\end{proposition}

\begin{theorem}[Cauchy--Schwarz]
\label{theorem-cauchy-schwarz}
For all $x,y\in H$,
$
|\langle x,y\rangle|\leq \|x\|\|y\|.
$
Equality holds if and only if $x,y$ are
linearly dependent.
\end{theorem}

\begin{proof}
Assume $y\neq0$. The vector
$
x-\frac{\langle x,y\rangle}
{\langle y,y\rangle}y
$
is orthogonal to $y$. By Pythagoras,
\begin{align*}
\|x\|^2
&=
\left\|
x-\frac{\langle x,y\rangle}
{\langle y,y\rangle}y
\right\|^2
+
\left\|
\frac{\langle x,y\rangle}
{\langle y,y\rangle}y
\right\|^2\\
&\geq
\frac{|\langle x,y\rangle|^2}{\|y\|^2}.
\end{align*}
Thus
$|\langle x,y\rangle|\leq\|x\|\|y\|$.
Equality holds exactly when the orthogonal
remainder is zero, i.e. when $x$ and $y$ are
linearly dependent.
\end{proof}

\begin{definition}
\label{definition-hilbert-space}
A {\it Hilbert space} is an inner product
space which is Banach for the norm induced by
the inner product.
\end{definition}

\begin{theorem}[Parallelogram law]
\label{theorem-parallelogram-law}
A norm on a vector space comes from an inner
product if and only if
$$
\|x+y\|^2+\|x-y\|^2=2\|x\|^2+2\|y\|^2
$$
for all $x,y$.
\end{theorem}

\begin{definition}
\label{definition-orthonormal-set}
A subset $S\subset H$ is {\it orthonormal} if
$\|s\|=1$ for all $s\in S$ and $\langle
s,t\rangle=0$ for $s\neq t$.
\end{definition}


\begin{proposition}
\label{proposition-maximal-orthonormal-extension}
Every orthonormal subset of an inner product
space is contained in a maximal orthonormal
subset.
\end{proposition}

\begin{proof}
Order the family of orthonormal subsets
containing the given set by inclusion. The
union of a chain is again orthonormal, so
Zorn's lemma gives a maximal element.
\end{proof}

\begin{remark}
\label{remark-maximal-orthonormal-characterization}
An orthonormal set $S$ is maximal if and only
if
$$
x\perp S\quad\Longrightarrow\quad x=0.
$$
\end{remark}

\begin{definition}
\label{definition-summable-family}
Let $(x_i)_{i\in I}$ be a family in a normed
space $X$. We say that it is {\it summable
to $x$} and write
$$
\sum_{i\in I}x_i=x
$$
if for every $\varepsilon>0$ there is a
finite $I_0\subseteq I$ such that, for every
finite $J\supseteq I_0$,
$$
\left\|
x-\sum_{i\in J}x_i
\right\|<\varepsilon.
$$
Equivalently, the net of finite partial sums
converges to $x$.
\end{definition}

\begin{proposition}
\label{proposition-summable-family-countable-support}
If $(x_i)_{i\in I}$ is summable, then
$$
\{i\in I:x_i\neq0\}
$$
is at most countable.
\end{proposition}

\begin{proof}
For every $\varepsilon>0$, the set
$$
I_\varepsilon=
\{i\in I:\|x_i\|>\varepsilon\}
$$
is finite. Indeed, if
$x=\sum_{i\in I}x_i$, choose finite
$I_0\subseteq I$ such that all finite partial
sums containing $I_0$ are within
$\varepsilon/2$ of $x$. If $i\notin I_0$,
then
\begin{align*}
\|x_i\|
&\leq
\left\|
x-\sum_{j\in I_0\cup\{i\}}x_j
\right\|
+
\left\|
x-\sum_{j\in I_0}x_j
\right\|\\
&<\varepsilon.
\end{align*}
Hence
$$
\{i:x_i\neq0\}
=
\bigcup_{n=1}^\infty
\{i:\|x_i\|>1/n\}
$$
is countable.
\end{proof}

\begin{theorem}[Bessel inequality]
\label{theorem-bessel}
If $S\subset H$ is orthonormal, then for
every $x\in H$,
$
\sum_{s\in S}|\langle x,s\rangle|^2\leq
\|x\|^2.
$
\end{theorem}

\begin{proof}
It is enough to prove that for every finite
subset $F\subset S$,
$
\sum_{s\in F}|\langle x,s\rangle|^2
\leq \|x\|^2.
$
Indeed, for a nonnegative family
$(a_s)_{s\in S}$, if every finite partial sum
is bounded above by $A$, define
$
a=\sup_{F\subset S,\ F\text{ finite}}
\sum_{s\in F}a_s\leq A.
$
Given $\varepsilon>0$, choose a finite
$F_0\subset S$ such that
$
a-\varepsilon<
\sum_{s\in F_0}a_s\leq a.
$
If $F\supset F_0$ is finite, then
$
a-\varepsilon<
\sum_{s\in F_0}a_s
\leq
\sum_{s\in F}a_s
\leq a,
$
so the finite partial sums converge to $a$.

Now fix a finite $F\subset S$. Since norms
are nonnegative,
$
0\leq
\left\|
x-\sum_{s\in F}\langle x,s\rangle s
\right\|^2.
$
Expanding the inner product and using
orthonormality,
\begin{align*}
\left\|
x-\sum_{s\in F}\langle x,s\rangle s
\right\|^2
&=
\|x\|^2
-2\sum_{s\in F}|\langle x,s\rangle|^2
+\sum_{s\in F}|\langle x,s\rangle|^2\\
&=
\|x\|^2-
\sum_{s\in F}|\langle x,s\rangle|^2.
\end{align*}
Therefore
$
\sum_{s\in F}|\langle x,s\rangle|^2
\leq \|x\|^2
$
for every finite $F\subset S$, and hence
$
\sum_{s\in S}|\langle x,s\rangle|^2
\leq \|x\|^2.
$
\end{proof}

\begin{theorem}[Parseval]
\label{theorem-parseval}
Let $H$ be a Hilbert space and let $S\subset
H$ be orthonormal. The following are
equivalent.
\begin{enumerate}
\item $S$ is maximal orthonormal.
\item The closed linear span of $S$ is $H$.
\item For every $x\in H$,
$
x=\sum_{s\in S}\langle x,s\rangle s.
$
\item For every $x\in H$,
$
\|x\|^2=\sum_{s\in S}|\langle x,s\rangle|^2.
$
\end{enumerate}
\end{theorem}

\begin{proof}
First, the series
$
\sum_{s\in S}\langle x,s\rangle s
$
exists for every $x\in H$. By the Cauchy
criterion, it is enough to show that for
every $\varepsilon>0$ there is a finite
$S_0\subset S$ such that for every finite
$S_1\subset S\setminus S_0$,
$
\left\|\sum_{s\in S_1}
\langle x,s\rangle s\right\|<\varepsilon.
$
By Pythagoras,
$
\left\|\sum_{s\in S_1}
\langle x,s\rangle s\right\|^2
=
\sum_{s\in S_1}|\langle x,s\rangle|^2.
$
By Bessel's inequality, the family
$(|\langle x,s\rangle|^2)_{s\in S}$ is
summable, so the Cauchy criterion holds.

Suppose now that $S$ is maximal. We have
$
x-\sum_{s\in S}\langle x,s\rangle s
\perp S.
$
Hence maximality gives
$
x=\sum_{s\in S}\langle x,s\rangle s.
$

If $S$ is not maximal, there exists
$y_0\in H$ with $\|y_0\|=1$ and
$y_0\perp S$. Hence
$
y_0\perp\overline{\operatorname{span}}(S),
$
so
$y_0\notin\overline{\operatorname{span}}(S)$.
This proves the equivalence with the
closed-span condition. It also gives
$
1=\|y_0\|^2
\neq
\sum_{s\in S}|\langle y_0,s\rangle|^2=0,
$
so Parseval's equality fails when $S$ is not
maximal.

Finally, if
$
x=\sum_{s\in S}\langle x,s\rangle s,
$
then Pythagoras gives
$
\|x\|^2=
\sum_{s\in S}|\langle x,s\rangle|^2.
$
\end{proof}



\begin{proposition}
\label{proposition-hilbert-is-l2-set}
Every Hilbert space is isometrically
isomorphic to $\ell^2(S)$ for some set $S$.
\end{proposition}

\begin{theorem}[Riesz representation]
\label{theorem-riesz-representation-hilbert}
Let $H$ be a Hilbert space. For every $f\in
H^*$ there is a unique $y\in H$ such that
$
f(x)=\langle x,y\rangle
$
for all $x\in H$. Moreover $\|f\|=\|y\|$.
\end{theorem}

\begin{proof}
Let $S\subseteq H$ be a maximal orthonormal
set. Formally define
$
y=
\sum_{s\in S}\overline{f(s)}\,s.
$
We first check that this family is summable.
For every finite $F\subseteq S$,
\begin{align*}
\|f\|
&\geq
\left|
f\left(
\frac{\sum_{s\in F}\overline{f(s)}s}
{\left\|\sum_{s\in F}\overline{f(s)}s\right\|}
\right)
\right|\\
&=
\left(
\sum_{s\in F}|f(s)|^2
\right)^{1/2}.
\end{align*}
Thus the nonnegative family
$(|f(s)|^2)_{s\in S}$ is summable, and
Pythagoras gives convergence of the vector
series defining $y$.

For $z\in H$, Parseval gives
$
z=\sum_{s\in S}\langle z,s\rangle s.
$
Therefore
\begin{align*}
\langle z,y\rangle
&=
\sum_{s\in S}
f(s)\langle z,s\rangle\\
&=
f\left(
\sum_{s\in S}\langle z,s\rangle s
\right)
=f(z).
\end{align*}
Cauchy--Schwarz gives $\|f\|\leq\|y\|$, and
evaluating at $y/\|y\|$ gives the reverse
inequality when $y\neq0$. Uniqueness follows
because
$
\langle z,y_1-y_2\rangle=0
\quad\text{for every }z
$
implies $y_1=y_2$.
\end{proof}

\begin{corollary}
\label{corollary-hilbert-dual-hilbert}
The dual of a Hilbert space is again a
Hilbert space.
\end{corollary}

\begin{corollary}
\label{corollary-hilbert-reflexive}
Every Hilbert space is reflexive.
\end{corollary}

\begin{proof}
The Riesz isometry identifies
$\overline H$ with $H^*$ and
$\overline{H^*}$ with $H^{**}$. Under these
identifications the canonical embedding
$J:H\to H^{**}$ is onto: if
$f=\langle\cdot,y\rangle_H$, then the
functional corresponding to $x\in H$
evaluates $f$ as
$
\langle x,y\rangle_H=f(x)=Jx(f).
$
\end{proof}

\begin{theorem}[Projection theorem]
\label{theorem-projection-hilbert}
Let $H$ be a Hilbert space and let $M\subset
H$ be a closed linear subspace. For every
$x\in H$ there is a unique $m\in M$ such that
$
\|x-m\|=d(x,M).
$
Moreover $x-m\perp M$.
\end{theorem}

\begin{proof}
Existence follows from the reflexive-space
minimization argument in the course notes:
for every $\varepsilon>d(x,M)$, the set
$
\overline{B(x,\varepsilon)}\cap M
$
is nonempty and weakly compact. These sets
are nested, so compactness gives a point
$m\in M$ with
$\|x-m\|=d(x,M)$.

Uniqueness follows from strict convexity of a
Hilbert space. If $m\neq m'$ both realize the
distance, the parallelogram law gives
$
\left\|
x-\frac{m+m'}2
\right\|
<d(x,M),
$
a contradiction.

Finally, suppose $h\in M$ and
$\langle x-m,h\rangle\neq0$. Choosing a
scalar $\lambda$ of sufficiently small
magnitude and opposite phase to
$\langle x-m,h\rangle$ gives
$
\|x-m+\lambda h\|^2
=
\|x-m\|^2
+2\operatorname{Re}
\langle x-m,\lambda h\rangle
+|\lambda|^2\|h\|^2
<
\|x-m\|^2,
$
contradicting minimality. Hence
$x-m\perp M$.
\end{proof}

\begin{definition}
\label{definition-orthogonal-projection}
The map $P_M:H\to M$ sending $x$ to the
unique nearest point of $M$ is called the
{\it orthogonal projection} onto $M$.
\end{definition}

\begin{proposition}
\label{proposition-orthogonal-projection}
The orthogonal projection $P_M$ is linear,
bounded, satisfies $P_M^2=P_M$, and has norm
$1$ if $M\neq0$.
\end{proposition}

\begin{proposition}
\label{example-fourier-orthonormal-system}
\begin{enumerate}
\item In $L_2(-\pi,\pi)$, the collection (or any subset thereof)
\[x_n=\frac{1}{\sqrt{2\pi}}e^{int},\qquad n=0,\pm1,\ldots\]
is an orthonormal set of vectors.
\begin{proof}
For any $n\in\mathbb{Z}$,
\[\int_{-\pi}^\pi x_n\overline{x}_ndt=\frac{1}{2\pi}\int_{-\pi}^\pi e^{int}\overline{e^{int}}dt=\frac{1}{2\pi}\int_{-\pi}^\pi \Vert e^{int}\Vert^2dt=1,\]
and if $m$ is another integer,
\[\int_{-\pi}^\pi x_n\overline{x}_mdt=\frac{1}{2\pi}\int_{-\pi}^\pi e^{int}\overline{e^{imt}}dt=\frac{1}{2\pi}\int_{-\pi}^\pi e^{(n-m)it}dt=\frac{1}{2\pi}\left[\frac{e^{(n-m)it}}{(n-m)i}\right]_{-\pi}^\pi=0.\]
\end{proof}
\item If we restric out attention to only real-valued functions that are square-integrable on the interval $[-\pi,\pi]$, then the collection (or any subset thereof)
\begin{align*}
\frac{1}{\sqrt{2\pi}},&\frac{1}{\sqrt{\pi}}\cos t,\frac{1}{\sqrt{\pi}}\cos 2t,\ldots\\
&\frac{1}{\sqrt{\pi}}\sin t,\frac{1}{\sqrt{\pi}}\sin 2t,\ldots
\end{align*}
is an orthonormal set.
\end{enumerate}
\end{proposition}

\begin{proposition}
\label{proposition-fourier-system-maximal}
The orthonormal system
$
S=
\left\{
\frac{e^{int}}{\sqrt{2\pi}}:n\in\mathbb Z
\right\}
$
is maximal in $L_2[0,2\pi]$.
\end{proposition}

\begin{proof}
We only prove maximality. Let
$f\in L_2[0,2\pi]$ satisfy
$
f\perp S.
$
We show that $f=0$.

First we replace $f$ by a continuous
function. Define
$
G(x)=\int_0^x f\,d\mu.
$
Then $G$ is absolutely continuous and
$G'=f$ almost everywhere. Hence, for every
$k\in\mathbb C$,
$
(G+k)'\perp S.
$
For $n\neq0$, integration by parts gives
$
\begin{aligned}
0
&=\int_0^{2\pi}(G+k)'e^{int}\,d\mu\\
&=G(2\pi)-G(0)
 -\int_0^{2\pi}(G+k)(in)e^{int}\,d\mu\\
&=-\int_0^{2\pi}(G+k)(in)e^{int}\,d\mu,
\end{aligned}
$
where $G(2\pi)-G(0)=\int_0^{2\pi}f\,d\mu=0$
because $f$ is orthogonal to the constant
function. Therefore
$
\int_0^{2\pi}(G+k)e^{int}\,d\mu=0
\qquad
\text{for every }n\in\mathbb Z\setminus
\{0\}.
$

Choose $k\in\mathbb C$ so that
$
H:=G+k
$
is also orthogonal to the constant function.
Then $H\perp S$.

Since $H$ is continuous, the
Stone--Weierstrass theorem gives, for every
$\varepsilon>0$, a trigonometric polynomial
$T\in\operatorname{span}(S)$ such that
$
\sup_{t\in[0,2\pi]}|H(t)-T(t)|
<\varepsilon.
$
Thus
$
H\in\overline{\operatorname{span}(S)}.
$
But $H\perp S$, hence
$H\perp\overline{\operatorname{span}(S)}$.
Therefore $H=0$. Since
$
f=G'=H'
$
almost everywhere, we conclude that $f=0$.
Thus $S$ is maximal.
\end{proof}


\section{Schauder bases}
\label{section-schauder-bases}

\begin{definition}
\label{definition-schauder-basis}
Let $X$ be a Banach space. A sequence
$(x_n)\subseteq X$ is a {\it Schauder basis}
of $X$ if for every $x\in X$ there exists a
unique scalar sequence $(\alpha_n)$ such that
$
x=\sum_{n=1}^\infty\alpha_nx_n.
$
\end{definition}

\begin{example}
\label{example-schauder-bases}
\begin{enumerate}
\item The canonical sequence $(e_n)$ is a
Schauder basis of $\ell_p$ for
$1\leq p<\infty$, and of $c_0$.
\item The space $\ell_\infty$ has no Schauder
basis, since every Banach space with a
Schauder basis is separable.
\item The course notes record Enflo's 1973
example of a separable Banach space without a
Schauder basis.
\end{enumerate}
\end{example}

\begin{theorem}
\label{theorem-schauder-basis-coordinate-functionals}
Let $X$ be a Banach space and
$(x_n)\subseteq X$. The following are
equivalent.
\begin{enumerate}
\item $(x_n)$ is a Schauder basis of $X$.
\item There exist functionals
$(x_n^*)\subseteq X^*$ such that
$
x_n^*(x_m)=\delta_{nm}
$
and
$
x=\sum_{n=1}^\infty x_n^*(x)x_n
\qquad\text{for every }x\in X.
$
\end{enumerate}
\end{theorem}

\begin{proof}
Assume (2). If
$
\sum_n\alpha_nx_n=\sum_n\beta_nx_n,
$
applying $x_k^*$ gives
$\alpha_k=\beta_k$ for every $k$, so the
representation is unique.

Now assume (1). Uniqueness gives algebraic
coordinate functionals $x_n^*$ satisfying
$
x=\sum_nx_n^*(x)x_n.
$
Let
$
S_Nx=\sum_{n=1}^Nx_n^*(x)x_n.
$
Since
$
x_n^*(x)x_n=S_nx-S_{n-1}x,
$
it is enough to show every $S_N$ is
continuous.

Following the course notes, define
$
\|x\|'=\sup_N\|S_Nx\|.
$
The supremum is finite because $S_Nx\to x$,
and
$
\|x\|\leq\|x\|'.
$
The notes prove that $(X,\|\cdot\|')$ is
complete by taking a $\|\cdot\|'$-Cauchy
sequence $(y_m)$, first taking its limit in
the original norm, and then using the
finite-dimensional ranges of the $S_N$ to
show convergence in $\|\cdot\|'$ as well.
Consequently the identity map
$
(X,\|\cdot\|')\longrightarrow(X,\|\cdot\|)
$
is a continuous bijection between Banach
spaces. By the bounded inverse theorem, the
two norms are equivalent. Since
$
\|S_Nx\|\leq\|x\|',
$
each $S_N$ is continuous in the original
norm, and hence so is each $x_n^*$.

The Portuguese source contains several index
typos inside the Cauchy-sequence computation;
the argument above records its stated
construction and conclusion.
\end{proof}

\begin{definition}
\label{definition-basis-constant}
For a Schauder basis $(x_n)$ with partial-sum
operators $S_n$, the number
$
K=\sup_n\|S_n\|
$
is its {\it basis constant}. The basis is
{\it monotone} if $K=1$.
\end{definition}

\begin{definition}
\label{definition-basic-sequence}
A sequence $(x_n)$ in a Banach space $X$ is
{\it basic} if it is a Schauder basis of
$
\overline{\operatorname{span}}\{x_n:n\in
\mathbb N\}.
$
\end{definition}

\begin{proposition}
\label{proposition-basic-sequence-characterization}
A sequence
$(x_n)\subseteq X\setminus\{0\}$ is basic if
and only if there is $L\geq1$ such that
$
\left\|
\sum_{n=1}^m\alpha_nx_n
\right\|
\leq
L
\left\|
\sum_{n=1}^k\alpha_nx_n
\right\|
$
for every $k\geq m$ and every choice of
scalars $\alpha_1,\ldots,\alpha_k$.
\end{proposition}

\begin{proof}
If $(x_n)$ is basic, the inequality follows
from uniform boundedness of the partial-sum
operators.

Conversely, let
$
Y=\operatorname{span}\{x_n:n\in\mathbb N\}
$
and define on $Y$
$
S_m\left(\sum_{n=1}^k\alpha_nx_n\right)
=
\sum_{n=1}^{\min\{m,k\}}\alpha_nx_n.
$
The hypothesis gives $\|S_m\|\leq L$, so each
$S_m$ extends continuously to $\overline Y$.
Moreover,
$
S_mS_k=S_kS_m=S_{\min\{m,k\}},
$
the range of $S_m$ is
$\operatorname{span}\{x_1,\ldots,x_m\}$, and
$S_mx\to x$ first on $Y$ and then, by the
uniform bound, on $\overline Y$. These facts
give the unique Schauder expansion on
$\overline Y$.
\end{proof}

\begin{definition}
\label{definition-biorthogonal-functionals}
For a basic sequence $(x_n)$, the coordinate
functionals $(x_n^*)$ are called its
{\it biorthogonal functionals}.
\end{definition}

\begin{theorem}
\label{theorem-equivalent-basic-sequences}
Let $(x_n)\subset X$ and $(y_n)\subset Y$ be
basic sequences. The following are
equivalent.
\begin{enumerate}
\item For every scalar sequence $(\alpha_n)$,
the series $\sum\alpha_nx_n$ converges if and
only if $\sum\alpha_ny_n$ converges.
\item There is an isomorphism
$
T:
\overline{\operatorname{span}}\{x_n\}
\longrightarrow
\overline{\operatorname{span}}\{y_n\}
$
such that $Tx_n=y_n$ for every $n$.
\end{enumerate}
\end{theorem}

\begin{proof}
The implication (2)$\Rightarrow$(1) is
immediate. Assume (1). Define $T$ by
$
T\left(\sum_n\alpha_nx_n\right)
=
\sum_n\alpha_ny_n.
$
Condition (1) makes $T$ a well-defined
bijection. To prove continuity, the course
notes use the closed graph theorem. If
$z_m\to z$ and $Tz_m\to w$, then for every
$j$,
$
y_j^*(w)
=
\lim_m y_j^*(Tz_m)
=
\lim_m x_j^*(z_m)
=
x_j^*(z)
=
y_j^*(Tz).
$
Uniqueness of Schauder coefficients gives
$w=Tz$. Thus $T$ has closed graph and is
bounded; the same argument applies to
$T^{-1}$.
\end{proof}

\begin{corollary}
\label{corollary-equivalent-basic-sequence-estimate}
Under the equivalent conditions above, there
is $L\geq1$ such that for every $n$ and every
$\alpha_1,\ldots,\alpha_n$,
$
\frac1L
\left\|\sum_{i=1}^n\alpha_ix_i\right\|
\leq
\left\|\sum_{i=1}^n\alpha_iy_i\right\|
\leq
L
\left\|\sum_{i=1}^n\alpha_ix_i\right\|.
$
\end{corollary}

\subsection{Constructing basic sequences}
\label{subsection-constructing-basic-sequences}

\begin{theorem}
\label{theorem-weak-closure-basic-sequence}
Let $S\subseteq X$ satisfy
$
0\notin\overline S^{\,\|\cdot\|}
\qquad\text{and}\qquad
0\in\overline S^{\,w}.
$
Then $S$ contains a basic sequence.
\end{theorem}

\begin{lemma}
\label{lemma-basic-sequence-finite-dimensional-step}
Let $E\subseteq X^*$ be finite-dimensional
and let
$S\subseteq X^*$ satisfy
$
0\notin\overline S^{\,\|\cdot\|},
\qquad
0\in\overline S^{\,w}.
$
For every $\varepsilon>0$ there exists
$x^*\in S$ such that
$
\|e^*+\lambda x^*\|
\geq
(1-\varepsilon)\|e^*\|
$
for every $e^*\in E$ and every real
$\lambda$.
\end{lemma}

\begin{proof}
Fix $\varepsilon>0$. Choose $\delta>0$.
Since $S_E$ is compact, choose
$e_1^*,\ldots,e_N^*\in S_E$ forming a
$\delta$-net. For every $i$, choose
$e_i\in S_X$ with
$
e_i^*(e_i)>1-\delta.
$
Since $0\in\overline S^{\,w}$, choose
$x^*\in S$ such that
$
|x^*(e_i)|<\delta
\qquad(i=1,\ldots,N).
$
Set
$
\alpha=\inf_{y^*\in\overline S}\|y^*\|>0.
$
For $e^*\in S_E$, choose $e_i^*$ within
$\delta$ of $e^*$.

The source splits into the two cases
$|\lambda|>2/\alpha$ and
$|\lambda|\leq2/\alpha$; the second condition
is printed there a second time as
$|\lambda|>2/\alpha$. In the first case,
$
\|e^*+\lambda x^*\|
\geq|\lambda|\,\|x^*\|-1\geq1.
$
In the second,
\begin{align*}
\|e^*+\lambda x^*\|
&\geq
|(e^*+\lambda x^*)(e_i)|\\
&\geq
1-2\delta-\frac{2\delta}{\alpha}.
\end{align*}
Choosing $\delta$ so that the last expression
is at least $1-\varepsilon$ proves the lemma;
homogeneity gives the general $e^*\in E$.
\end{proof}

\begin{proof}[Proof of Theorem
\ref{theorem-weak-closure-basic-sequence}]
View $X$ canonically in $X^{**}$. Choose
positive numbers $(\delta_n)$ such that
$
\prod_n(1-\delta_n)>\frac1{1+\varepsilon}.
$
Iterating Lemma
\ref{lemma-basic-sequence-finite-dimensional-step}
produces $(x_n)\subseteq S$ such that for
every $m$ and scalars
$a_1,\ldots,a_{m+1}$,
$
(1-\delta_m)
\left\|
\sum_{n=1}^m a_nx_n
\right\|
\leq
\left\|
\sum_{n=1}^{m+1}a_nx_n
\right\|.
$
Iterating this estimate gives, for
$m\leq k$,
$
\left\|
\sum_{n=1}^m a_nx_n
\right\|
\leq
(1+\varepsilon)
\left\|
\sum_{n=1}^k a_nx_n
\right\|.
$
Proposition
\ref{proposition-basic-sequence-characterization}
shows that $(x_n)$ is basic.

The Portuguese source has several index
typos in this iteration; the displayed
product estimate is the one it explicitly
states as the goal of the construction.
\end{proof}

\begin{corollary}
\label{corollary-weak-null-basic-subsequence}
If $(x_n)\subseteq X$ converges weakly to
zero and
$
\inf_n\|x_n\|>0,
$
then $(x_n)$ has a basic subsequence.
\end{corollary}

\begin{theorem}
\label{theorem-infinite-dimensional-basic-sequence}
Every infinite-dimensional Banach space
contains a basic sequence.
\end{theorem}

\begin{proof}
Apply Theorem
\ref{theorem-weak-closure-basic-sequence}
to $S=S_X$. Clearly
$0\notin\overline{S_X}^{\,\|\cdot\|}$.
If $0\notin\overline{S_X}^{\,w}$, there
would be
$f_1,\ldots,f_n\in X^*$ and
$\varepsilon>0$ such that
$
\{x:|f_i(x)|<\varepsilon\text{ for all }i\}
\cap S_X
=\varnothing.
$
But in an infinite-dimensional space
$
\bigcap_{i=1}^n\ker f_i
$
contains a nonzero vector; after
normalizing it we obtain a point of the
displayed intersection, a contradiction.
\end{proof}

\begin{definition}
\label{definition-vector-valued-lp-c0}
Let $X$ be a Banach space. For
$1\leq p<\infty$, define
$
\ell_p(X)
=
\left\{(x_n)\in X^{\mathbb N}:
\sum_n\|x_n\|^p<\infty
\right\}
$
with
$
\|(x_n)\|
=
\left(\sum_n\|x_n\|^p\right)^{1/p}.
$
Also define
$
c_0(X)
=
\left\{(x_n)\in\ell_\infty(X):
\|x_n\|\to0
\right\}
$
with the supremum norm.
\end{definition}

\section{Banach algebras}
\label{section-banach-algebras}

\begin{definition}
\label{definition-banach-algebra}
A {\it Banach algebra} $A$ is a Banach space
equipped with an associative multiplication
such that $\|ab\| \leq \|a\| \|b\|$ for all
$a, b \in A$. If there exists $1 \in A$ such
that $a1 = 1a = a$, $A$ is a {\it unital
Banach algebra}.
\end{definition}

\begin{definition}
\label{definition-normed-algebra-ideal-homomorphism}
Let $A$ be an algebra.
\begin{enumerate}
\item A vector subspace $I\subseteq A$ is a
{\it right ideal} if $ab\in I$ whenever
$a\in A$ and $b\in I$.
\item It is a {\it left ideal} if $ab\in I$
whenever $a\in I$ and $b\in A$.
\item It is a {\it two-sided ideal} if both
conditions hold.
\item A linear map $T:A\to B$ between
algebras is a {\it homomorphism} if
$
T(ab)=T(a)T(b).
$
\end{enumerate}
\end{definition}

\begin{example}
\label{example-banach-algebra-ideals}
\begin{enumerate}
\item For $x_0\in K$,
$
I_{x_0}=\{f\in C(K):f(x_0)=0\}
$
is an ideal of $C(K)$.
\item $K(X)$ is an ideal of
$\mathcal L(X)$.
\item $c_0$ is an ideal of $\ell_\infty$.
\item The evaluation map
$f\mapsto f(x_0)$ on $C(K)$ is an algebra
homomorphism.
\end{enumerate}
Kernels of algebra homomorphisms are ideals.
\end{example}

\begin{example}
\label{example-banach-algebras}
\begin{enumerate}
\item $C(K)$ for a compact Hausdorff space
$K$ with point-wise multiplication.
\item $\mathcal{L}(X)$ for a Banach space $X$
with composition.
\item $L^1(\mathbb{R})$ with the convolution
product $(f \ast g)(t) = \int f(t-s)g(s) ds$.
\item $K(X)$, the space of compact operators,
which is a closed ideal in $\mathcal{L}(X)$.
\end{enumerate}
\end{example}

\begin{proposition}
\label{proposition-algebra-quotient}
If $I \subseteq A$ is a closed two-sided
ideal, then $A/I$ is a Banach algebra with
the quotient norm.
\end{proposition}

\begin{proof}
The product is
$
(a+I)(b+I)=ab+I.
$
It is well defined: if $a-a'\in I$ and
$b-b'\in I$, then
$
ab-a'b'
=
a(b-b')+(a-a')b'
\in I.
$
For the quotient norm, for arbitrary
$a',b'\in I$,
$
ab+ab'+a'b+a'b'=(a+a')(b+b'),
$
so
\begin{align*}
\|ab+I\|
&\leq
\inf_{a',b'\in I}
\|(a+a')(b+b')\|\\
&\leq
\left(\inf_{a'\in I}\|a+a'\|\right)
\left(\inf_{b'\in I}\|b+b'\|\right)\\
&=
\|a+I\|\,\|b+I\|.
\end{align*}
Completeness follows from completeness of
the quotient of a Banach space by a closed
subspace.
\end{proof}

\begin{definition}
\label{definition-spectrum-algebra}
Let $A$ be a unital Banach algebra and $a \in
A$. The {\it spectrum} of $a$ is
$
\sigma(a) = \{ \lambda \in \mathbb{C} : a -
\lambda 1 \text{ is not invertible in } A \}.
$
The {\it spectral radius} of $a$ is $r(a) =
\sup \{ |\lambda| : \lambda \in \sigma(a)
\}$.
\end{definition}

\begin{proposition}[Polynomial spectral mapping]
\label{proposition-polynomial-spectral-mapping}
Let $a\in A$ and let $p\in\mathbb C[z]$. Then
$
p(\sigma(a))=\sigma(p(a)).
$
\end{proposition}

\begin{proof}
The constant case is immediate. Otherwise,
for $\lambda\in\mathbb C$ factor
$
p(z)-\lambda
=
c(z-\lambda_1)\cdots(z-\lambda_n),
\qquad c\neq0.
$
Then
$
p(a)-\lambda
=
c(a-\lambda_1)\cdots(a-\lambda_n).
$
The factors commute. A product of commuting
elements is invertible if and only if every
factor is invertible. Hence
$
\lambda\in\sigma(p(a))
\iff
\lambda_i\in\sigma(a)
\text{ for some }i
\iff
\lambda\in p(\sigma(a)).
$
\end{proof}

\begin{proposition}[Neumann series]
\label{proposition-neumann-series-banach-algebra}
If $\|a\|<1$, then $1-a$ is invertible and
$
(1-a)^{-1}=\sum_{n=0}^\infty a^n.
$
\end{proposition}

\begin{proof}
The series converges because
$
\sum_{n=0}^\infty\|a^n\|
\leq
\sum_{n=0}^\infty\|a\|^n<\infty.
$
If
$b=\sum_{n=0}^\infty a^n$, then
\begin{align*}
(1-a)b
&=
\lim_{N\to\infty}
(1-a)(1+a+\cdots+a^N)\\
&=
\lim_{N\to\infty}(1-a^{N+1})
=1,
\end{align*}
and similarly $b(1-a)=1$.
\end{proof}

\begin{corollary}
\label{corollary-invertibles-open-banach-algebra}
The group of invertible elements
$\operatorname{Inv}(A)$ is open in $A$.
\end{corollary}

\begin{proof}
If $a$ is invertible and
$
\|a-b\|<\|a^{-1}\|^{-1},
$
then
$
\|1-ba^{-1}\|
=
\|(a-b)a^{-1}\|<1.
$
Thus $ba^{-1}$ is invertible by the Neumann
series, and therefore so is $b$.
\end{proof}

\begin{theorem}
\label{theorem-inversion-differentiable-banach-algebra}
The inversion map
$
\operatorname{Inv}(A)\longrightarrow
\operatorname{Inv}(A),
\qquad a\longmapsto a^{-1},
$
is differentiable, with
$
D(a\mapsto a^{-1})_a(b)
=
-a^{-1}ba^{-1}.
$
\end{theorem}

\begin{proof}
For $\|a^{-1}b\|<1/2$,
\begin{align*}
(a+b)^{-1}
&=
(1+a^{-1}b)^{-1}a^{-1}\\
&=
\left(
1-a^{-1}b+
\sum_{n=2}^\infty(-a^{-1}b)^n
\right)a^{-1}.
\end{align*}
Hence
\begin{align*}
\|(a+b)^{-1}-a^{-1}+a^{-1}ba^{-1}\|
&\leq
\sum_{n=2}^\infty
\|a^{-1}b\|^n\|a^{-1}\|\\
&\leq
2\|a^{-1}\|^3\|b\|^2.
\end{align*}
Dividing by $\|b\|$ and letting $b\to0$
proves the formula.

The Portuguese notes state the derivative
once with the opposite sign, but their
remainder computation is exactly the one
above and therefore yields the negative sign.
\end{proof}

\begin{theorem}[Gelfand]
\label{theorem-spectrum-nonempty}
For any $a \in A$, $\sigma(a)$ is a non-empty
compact subset of $\mathbb{C}$. Moreover,
$\sigma(a) \subseteq \{ z \in \mathbb{C} :
|z| \leq \|a\| \}$.
\end{theorem}

\begin{proof}
If $|\lambda|>\|a\|$, then
$\|\lambda^{-1}a\|<1$, so
$1-\lambda^{-1}a$ is invertible by the
Neumann series. Thus $a-\lambda$ is
invertible and
$
\sigma(a)\subseteq
\{\lambda:|\lambda|\leq\|a\|\}.
$
Since $\operatorname{Inv}(A)$ is open,
$\mathbb C\setminus\sigma(a)$ is open, so
$\sigma(a)$ is closed and therefore compact.

Suppose $\sigma(a)=\varnothing$. Then
$
\Phi(\lambda)=(a-\lambda)^{-1}
$
is defined on all of $\mathbb C$. For every
$f\in A^*$, the scalar function
$f\circ\Phi$ is entire. For
$|\lambda|>2\|a\|$,
\begin{align*}
\|(a-\lambda)^{-1}\|
&=
|\lambda|^{-1}
\|(1-\lambda^{-1}a)^{-1}\|\\
&\leq
\frac{2}{|\lambda|}.
\end{align*}
On the remaining compact disk, continuity of
the inverse map makes $\Phi$ bounded. Thus
$f\circ\Phi$ is bounded and hence constant
by Liouville. Since this holds for every
$f\in A^*$, Hahn--Banach forces
$\Phi$ itself to be constant, which is
impossible. Therefore $\sigma(a)$ is
nonempty.
\end{proof}

\begin{corollary}
\label{corollary-complex-banach-division-algebra}
If $A$ is a complex unital Banach algebra and
every nonzero element is invertible, then
$
A\cong\mathbb C.
$
\end{corollary}

\begin{proof}
For $a\in A$, choose
$\lambda\in\sigma(a)$. Then
$a-\lambda1$ is not invertible, so by
hypothesis it is zero. Thus every
$a\in A$ is a scalar multiple of the unit.
\end{proof}

\begin{theorem}[Spectral radius formula]
\label{theorem-spectral-radius-formula}
For any $a \in A$, $r(a) = \lim_{n \to
\infty} \|a^n\|^{1/n}$.
\end{theorem}

\begin{remark}
\label{remark-spectral-radius-source-proof-gap}
The Portuguese course notes begin a proof by
first observing, from polynomial spectral
mapping, that
$
r(a)\leq\inf_n\|a^n\|^{1/n}.
$
For the reverse inequality they consider
the resolvent power series
$
(1-\lambda a)^{-1}
$
and use Banach--Steinhaus to bound
$(\lambda^na^n)_n$. The source has an
unresolved marked gap at the point where the
domain of this analytic power series is
identified, and later interchanges
$\lambda$ and $\lambda^{-1}$. We therefore
retain the theorem here but do not silently
replace the source's incomplete proof by a
different proof.
\end{remark}

\begin{example}
\label{example-spectral-radius-strictly-less-norm}
Let
$
A=
\{f\in C[0,1]:f\text{ is }C^1\}
$
with
$
\|f\|=\|f\|_\infty+\|f'\|_\infty.
$
For the function $z(t)=t$,
$
\|z^n\|=1+n,
$
so
$
r(z)=\lim_n(1+n)^{1/n}=1<2=\|z\|.
$
\end{example}

\begin{theorem}
\label{theorem-spectrum-closed-subalgebra}
Let $A$ be a unital Banach algebra and
$B\subseteq A$ a closed subalgebra containing
$1_A$. Then:
\begin{enumerate}
\item $\operatorname{Inv}(B)$ is both open
and closed in
$B\cap\operatorname{Inv}(A)$.
\item For $b\in B$,
$
\sigma_A(b)\subseteq\sigma_B(b),
\qquad
\partial\sigma_B(b)
\subseteq
\partial\sigma_A(b).
$
\item If
$\mathbb C\setminus\sigma_A(b)$ is connected,
then
$
\sigma_A(b)=\sigma_B(b).
$
\end{enumerate}
\end{theorem}

\begin{proof}
The first set is open because invertibles are
open in every Banach algebra. To prove it is
closed in
$B\cap\operatorname{Inv}(A)$, let
$b_n\in\operatorname{Inv}(B)$ with
$b_n\to b\in B\cap\operatorname{Inv}(A)$.
Continuity of inversion in $A$ gives
$b_n^{-1}\to b^{-1}$, and closedness of $B$
implies $b^{-1}\in B$.

Since
$\operatorname{Inv}(B)\subseteq
\operatorname{Inv}(A)$,
one gets
$\sigma_A(b)\subseteq\sigma_B(b)$.
If
$\lambda\in\partial\sigma_B(b)$, take
$\lambda_n\notin\sigma_B(b)$ with
$\lambda_n\to\lambda$. If
$\lambda\notin\sigma_A(b)$, then
$b-\lambda_n$ converges inside
$B\cap\operatorname{Inv}(A)$ to
$b-\lambda$, and part (1) would force
$b-\lambda\in\operatorname{Inv}(B)$, a
contradiction. This proves the boundary
inclusion.

Finally,
$\mathbb C\setminus\sigma_B(b)$ is both open
and closed in the connected set
$\mathbb C\setminus\sigma_A(b)$, so the two
resolvent sets coincide.
\end{proof}

\begin{example}
\label{example-disk-algebra-spectrum}
Let $A$ be the disk algebra and let $B$ be
its image under restriction to
$\partial\mathbb D$. The source records
$
\sigma_{C(\partial\mathbb D)}(z)
=
\partial\mathbb D,
\qquad
\sigma_A(z)=\sigma_B(z)=\mathbb D.
$
\end{example}

\subsection{Exponentials in Banach algebras}
\label{subsection-exponential-banach-algebra}

\begin{definition}
\label{definition-exponential-banach-algebra}
For $a$ in a unital Banach algebra, define
$
e^a=
\sum_{n=0}^\infty\frac{a^n}{n!}.
$
\end{definition}

\begin{theorem}
\label{theorem-exponential-banach-algebra}
Let $A$ be a unital Banach algebra.
\begin{enumerate}
\item If $f:\mathbb R\to A$ satisfies
$f(0)=1$ and
$f'(t)=af(t)$, then
$
f(t)=e^{ta}.
$
\item $e^a$ is invertible and
$
(e^a)^{-1}=e^{-a}.
$
\item If $ab=ba$, then
$
e^{a+b}=e^ae^b.
$
\end{enumerate}
\end{theorem}

\begin{proof}
The exponential series converges uniformly
on bounded $t$-intervals and may be
differentiated term by term, so
$g(t)=e^{ta}$ satisfies
$
g(0)=1,\qquad g'(t)=ag(t).
$
If $f$ is another solution, put
$
h(t)=e^{-ta}f(t).
$
Using the product rule and commutation of
$a$ with its powers,
$
h'(t)
=
-ae^{-ta}f(t)+e^{-ta}af(t)=0.
$
Thus $h(t)=1$ and
$f(t)=e^{ta}$. Taking $f(t)=e^{ta}$ and
$t=\pm1$ gives
$e^ae^{-a}=1$.

If $ab=ba$, set
$
g(t)=e^{ta}e^{tb}.
$
Then
$
g'(t)
=
(a+b)g(t),
\qquad
g(0)=1.
$
Part (1) gives
$g(t)=e^{t(a+b)}$, and $t=1$ yields the
claim.

The Portuguese source has several typos in
this proof (including writing the same
exponential twice in the last part); the
ODE argument above is the argument it
records.
\end{proof}

\section
{Fredholm operators and the Calkin algebra}
\label{section-fredholm-theory}

\begin{definition}
\label{definition-fredholm-operator}
A bounded linear operator $T \in
\mathcal{L}(X, Y)$ is a {\it Fredholm
operator} if $\dim \ker T < \infty$ and
$\operatorname{codim} \operatorname{Im} T <
\infty$. The {\it index} of $T$ is
$
\operatorname{ind}(T) = \dim \ker T -
\operatorname{codim} \operatorname{Im} T.
$
\end{definition}

\begin{proposition}
\label{proposition-composition-fredholm-index}
If $T\in\mathcal L(X,Y)$ and $S\in\mathcal L(Y,Z)$ are Fredholm, then
$ST$ is Fredholm and
$
\operatorname{ind}(ST)=\operatorname{ind}S+\operatorname{ind}T.
$
\end{proposition}

\begin{proof}
Choose closed subspaces $Y_1,Y_2,Y_3$ so that
$
\operatorname{Im}T=(\ker S\cap\operatorname{Im}T)\oplus Y_1,
$
$
\ker S=(\ker S\cap\operatorname{Im}T)\oplus Y_2,
$
and
$
Y=\operatorname{Im}T\oplus Y_2\oplus Y_3.
$
The map
$
\ker(ST)\longrightarrow\ker S\cap\operatorname{Im}T,\qquad x\longmapsto Tx
$
is surjective with kernel $\ker T$, so
$
\dim\ker(ST)=\dim\ker T+\dim(\ker S\cap\operatorname{Im}T).
$
Also
$
\operatorname{Im}S=S(Y_1)+S(Y_3),\qquad \operatorname{Im}(ST)=S(Y_1).
$
Counting the finite-dimensional pieces gives
$
\operatorname{ind}(ST)=\operatorname{ind}S+\operatorname{ind}T.
$
\end{proof}

\begin{corollary}
\label{corollary-index-power}
If $T$ is Fredholm, then
$
\operatorname{ind}(T^n)=n\,\operatorname{ind}T.
$
\end{corollary}


\begin{theorem}[Atkinson]
\label{theorem-atkinson-calkin}
Let $X$ be an infinite-dimensional Banach
space. Let $\pi : \mathcal{L}(X) \to
\mathcal{L}(X)/K(X)$ be the natural
projection into the Calkin algebra. Then $T
\in \mathcal{L}(X)$ is Fredholm if and only
if $\pi(T)$ is invertible in
$\mathcal{L}(X)/K(X)$.
\end{theorem}

\begin{lemma}
\label{lemma-fredholm-parametrix-course}
If $T\in\mathcal L(X)$ is Fredholm, there is
$S\in\mathcal L(X)$ such that
$
I-ST
\quad\text{and}\quad
I-TS
$
have finite rank.
\end{lemma}

\begin{proof}
Choose closed complements
$
X=X_1\oplus\ker T,
\qquad
X=\operatorname{Im}T\oplus X_2.
$
The restriction
$
T|_{X_1}:X_1\longrightarrow\operatorname{Im}T
$
is a bounded bijection between Banach spaces,
so its inverse is bounded. Define
$
S(y+z)=(T|_{X_1})^{-1}y,
\qquad
y\in\operatorname{Im}T,\ z\in X_2.
$
Then $ST$ is the projection onto $X_1$ along
$\ker T$, while $TS$ is the projection onto
$\operatorname{Im}T$ along $X_2$. Thus
$I-ST$ and $I-TS$ have finite-dimensional
ranges.

The Portuguese source prints $1=ST$ in one
line; its construction gives the finite-rank
remainder above.
\end{proof}

\begin{proof}[Proof of Theorem
\ref{theorem-atkinson-calkin}]
If $T$ is Fredholm, take $S$ from Lemma
\ref{lemma-fredholm-parametrix-course}. In
the Calkin algebra,
$
\pi(S)\pi(T)=\pi(I)=\pi(T)\pi(S),
$
so $\pi(T)$ is invertible.

Conversely, suppose $\pi(T)$ is invertible.
Choose $S\in\mathcal L(X)$ representing its
inverse. Then
$
K=I-ST,
\qquad
K'=I-TS
$
are compact. If $Tx=0$, then $Kx=x$, so
$
\ker T\subseteq\ker(K-I).
$
The eigenspace of the compact operator $K$
for the nonzero eigenvalue $1$ is
finite-dimensional, hence $\dim\ker T<\infty$.

Also
$
TS=I-K'.
$
The compact-operator result for the
nonzero scalar $1$ gives finite codimension
of $\operatorname{Im}(I-K')$. Since
$
\operatorname{Im}(I-K')
=
\operatorname{Im}(TS)
\subseteq\operatorname{Im}T,
$
the range of $T$ has finite codimension.
Thus $T$ is Fredholm.
\end{proof}


\begin{proposition}
\label{proposition-index-stability}
\begin{enumerate}
\item The set of Fredholm operators $\Phi(X)$
is open in $\mathcal{L}(X)$.
\item The index map $\operatorname{ind} :
\Phi(X) \to \mathbb{Z}$ is continuous (and
thus locally constant).
\item If $T$ is Fredholm and $S$ is compact,
then $T+S$ is Fredholm and
$\operatorname{ind}(T+S) =
\operatorname{ind}(T)$.
\end{enumerate}
\end{proposition}

\section
{Spectral theory for compact operators}
\label{section-spectral-compact}

\begin{theorem}
\label{theorem-spectral-compact-ops}
Let $X$ be a Banach space and $T \in
\mathcal{K}(X)$ a compact operator.
\begin{enumerate}
\item $\sigma(T)$ is at most countable and
has no limit point except possibly $0$.
\item Every $\lambda \in \sigma(T) \setminus
\{0\}$ is an eigenvalue of finite
multiplicity.
\item $\operatorname{ind}(T - \lambda I) = 0$
for all $\lambda \neq 0$.
\end{enumerate}
\end{theorem}

\begin{theorem}
[Spectral theorem for compact self-adjoint
operators]
\label{theorem-spectral-compact-selfadjoint}
Let $H$ be a Hilbert space and $T \in
\mathcal{K}(H)$ be self-adjoint ($T = T^*$).
There exists a maximal orthonormal set $\{
e_n \}$ of eigenvectors of $T$ with real
eigenvalues $\{ \lambda_n \}$. For any $x \in
H$:
$$
Tx = \sum_n \lambda_n \langle x, e_n \rangle
e_n.
$$
\end{theorem}

\section
{Special topics and advanced examples}
\label{section-advanced-topics}

\begin{example}[Fourier Transform]
\label{example-fourier-l1}
For $f \in L^1(\mathbb{R})$, the Fourier
transform is $\widehat{f}(\xi) =
\int_\mathbb{R} e^{-2\pi i \xi x} f(x) dx$.
The map $\mathcal{F} : L^1(\mathbb{R}) \to
C_0(\mathbb{R})$ is a bounded linear operator
with $\|\mathcal{F}\| \leq 1$. It is
injective but not surjective.
\end{example}

\begin{example}
[Weak Derivatives and Sobolev Spaces]
\label{example-sobolev-w1p}
Let $f \in L^1_{\operatorname{loc}}(0,1)$. A
function $g$ is the {\it weak derivative} of
$f$ if
$$
\int_0^1 f(x)\varphi'(x) dx = -\int_0^1
g(x)\varphi(x) dx \quad \forall \varphi \in
C_c^1(0,1).
$$
The Sobolev space $W^{1,p}(0,1) = \{ f \in
L^p(0,1) : f' \in L^p(0,1) \}$ is a Banach
space with norm $\|f\| = \|f\|_p + \|f'\|_p$.
It is reflexive for $1 < p < \infty$.
\end{example}

\begin{example}[Orlicz Spaces]
\label{example-orlicz-space}
Let $\varphi$ be a continuous, convex,
increasing function on $[0, \infty)$ with
$\varphi(0)=0$. The Orlicz space
$L^\varphi(\mathbb{R})$ consists of
measurable functions such that $\int
\varphi(|f|/C) dx < \infty$ for some $C > 0$.
The Luxemburg norm is $\|f\|_\varphi = \inf
\{ C > 0 : \int \varphi(|f|/C) dx \leq 1 \}$.
\end{example}

\section{Further exercises}
\label{section-functional-analysis-exercises}

\begin{exercise}
\label{exercise-dual-c0-l1}
Prove that $c_0^* \cong \ell_1$ and $\ell_1^*
\cong \ell_\infty$.
\end{exercise}

\begin{exercise}
\label{exercise-weak-vs-strong}
Show that in an infinite-dimensional space,
the weak topology is strictly coarser than
the norm topology. Find a sequence that
converges weakly but not in norm.
\end{exercise}

\begin{exercise}
\label{exercise-hilbert-projection}
Let $M$ be a closed subspace of a Hilbert
space $H$. Show that $H = M \oplus M^\perp$
and that the orthogonal projection $P_M$ is a
self-adjoint operator of norm 1.
\end{exercise}

\end{document}

